\documentclass[ecta,nameyear,draft]{econsocart}
\usepackage{amsmath,amsfonts,amssymb,amsthm,booktabs,array,longtable}
\usepackage[T1]{fontenc}
\usepackage{mathpazo,microtype,needspace,xurl}
\usepackage{xr-hyper}
\RequirePackage[colorlinks,citecolor=blue,linkcolor=blue,urlcolor=blue]{hyperref}
\startlocaldefs
\newtheorem{theorem}{Theorem}
\newtheorem{proposition}[theorem]{Proposition}
\newtheorem{lemma}[theorem]{Lemma}
\newtheorem{corollary}[theorem]{Corollary}
\newcommand{\E}{\mathbb E}
\newcommand{\R}{\mathbb R}
\endlocaldefs
\setlength{\emergencystretch}{3em}
\makeatletter\def\copyright@text{Prepared for referee review}\makeatother

\setcounter{secnumdepth}{2}

\externaldocument[][nocite]{build/supp}[supp.pdf]
\externaldocument[][nocite]{build/appendix}[appendix.pdf]
\begin{document}
\begin{frontmatter}
\title{Neural Bellman Operators}
\runtitle{Neural Bellman Operators}
\begin{aug}
\author[id=au1,addressref={add1}]{\fnms{Qian}~\snm{QI}\ead[label=e90ECTA]{qiqian@pku.edu.cn}}
\address[id=add1]{Peking University}
\end{aug}
\begin{abstract}
This paper develops Neural Bellman Operators for constructive policy evaluation and feasible improvement under the original economic law. Attained continuations retain action witnesses and an original-optimum error account. We extend qualification-to-cost transfer to complete multi-date policies by coupling shared-parameter action restrictions with their endogenous occupation flows. A joint state-context construction provides one neural circuit with uniform guarantees beyond the acquired catalogue, including adaptively generated contexts under explicit coverage conditions. Stored continuations and dated error allocation yield a separate horizon-complexity account. Prospective finite-menu selection and polyhedral price-fee regions connect qualification to economic decisions. A frozen independent investment evaluation identifies several policy-cost improvements and a noninferiority guarantee while retaining an unresolved fine-basis ordering. Nonlinear recursive evaluations and strategic deviations receive their appropriate order certificates. All preceding constructions, proofs, broader applications and adverse computational evidence remain part of the paper.
 The complete-pattern extension retains shared choices before relaxing occupations. A transport-qualified residual-span bound treats singular transition approximations under the original Bellman optimum and enters a coupled residual--transition accuracy budget. A frozen shared-proposal study retains its negative query-work result.
\end{abstract}
\begin{keyword}
\kwd{Dynamic programming}\kwd{Neural approximation}\kwd{Policy certification}
\kwd{Information acquisition}\kwd{Computational economics}
\end{keyword}
\end{frontmatter}
\section{Introduction}\label{sec:introduction65}
A continuation approximation matters through the economic decision it produces. Prediction error, original-optimum policy loss and verification work are different objects. This paper connects them through a constructive Neural Bellman Operator (NBO), its attained own-future continuation and the complete feasible controller implemented under the original law.

The operator fits each date against its own attained future and retains feasible action witnesses. Its affine--ReLU construction and primitive error budget do not require unrestricted nonconvex fitting to succeed. A complementary recursive implementation uses verified curvature and protected conventional refinement; trained proposals change optional query locations. Residual qualification permits an attained shallow continuation to replace recursive value queries under a stated error account. Affine calibration separates its represented derivative from a conservative coefficientwise penalty.

Our next question is how much optimization accuracy an economic decision needs. Rational switch barriers identify every acquired action possible throughout a qualification sublevel. Original-law conditional contrasts price those actions, instead of charging the full cost range on each unresolved cell. Exact thresholds determine the entire region of tolerances supporting an ordering, cost-equivalence or replacement conclusion. A verified primal--dual gap places a returned readout in that region without a fitting-success model or additional sample. Labels and the continuation policy remain fixed in that first-date comparison.

The prospective investment experiment compares both bases, both qualification rules and both observation resolutions. Calibrated qualification improves both ReLU and Bernstein controllers and reverses their earlier matched cost ordering. Independent arithmetic identifies original-law gains; actionwise transfer preserves some comparisons at nonzero optimization errors without identical actions. The full catalogue retains adverse and unresolved cases. Conditional economic choices require separate charges and deployment counts, not an invented conversion of runtime or simulation rows into welfare.

We now pass from a first-date comparison to complete policies. Changing an earlier action changes the distribution of later acquired states. The recursive transfer therefore couples actual occupation flows with exclusions imposed by a common parameter vector. A separate joint-domain construction produces one state-context neural circuit, and proves its approximation and feasible-policy guarantees throughout a compact economic family rather than only on its training catalogue. These results extend the original NBO construction; they do not replace it by a finite-menu benchmark.

\paragraph{Relation to existing methods.}
The construction draws on Lipschitz extensions and distance transforms \citep{mcshane1934,felzenszwalb2012}. Approximate policy iteration, safe improvement and convex dynamic programming have established analyses \citep{munos2003,pirotta2013,scherrer2015,yang2020convex}; learned warm starts and optimization-based control provide related numerical methods \citep{sakaue2023warm,sambharya2023,agrawal2020control}. A common verifier does not establish the incremental value of its neural proposal.

Control variates, telescoping evaluation and time-uniform concentration are established tools \citep{munos2006variance,jiang2016,howard2020time,karampatziakis2021}, as are minimax-regret decisions and shortest-path difference constraints \citep{manski2004,floyd1962}. The new link is from exact qualification barriers and original-law action effects to economic stopping regions for a constructed continuation, not priority for these underlying principles.

The same readout coefficients determine actions in every observation cell. Separate switch tests can therefore admit an alternative that never wins its full menu, or combinations of alternatives that no common readout can select. Section~\ref{sec:compatible84} retains both restrictions. An exact support dual certifies each winner event and selected joint events; a second, cell-assignment dual transfers the resulting compatibility information to original-law cost. This strengthens the existing NBO optimization-to-policy account without changing a fitted label, future controller or economic primitive.

The additional calculation is deterministic qualification of previously purchased evidence. It reports the incremental work of full-menu and shared-readout checks separately from construction, integration and policy-cost inference. The economic enclosure remains valid when a numerical solver is inconclusive, because only independently verified multipliers may exclude an event. Standard convex duality supplies the lower-bound principle \citep{boyd2004}; the contribution here is its connection to acquired NBO actions, common coefficients and the paid attained-policy cost event.

The multi-date result retains this distinction between a certificate and an attained optimizer, while adding endogenous occupations and whole-cell transition bounds. Its nonlinear-recursion and strategic-deviation counterparts use order comparison, not an invalid expectation identity. The horizon account explicitly distinguishes reusable stored continuations from repeated descendant solves. Finally, a fixed-menu population guarantee and a separate calibration/evaluation split distinguish prospective selection from an ex post catalogue winner. Uniform price-fee regions show which decisions survive uncertain computational valuations.

\paragraph{Complete policies, singular transitions and paid precision.}
The common-coefficient restriction is imposed on an entire contingent pattern before its occupation flow is relaxed. This strengthens pairwise compatibility and supplies an exact finite-state limit with a separate continuous-law enclosure error. When transition-probability widths do not shrink, a complementary transport certificate tests the numerical kernel against the stored continuation itself. Its residual-span account bounds complete-policy loss relative to the original Bellman optimum even for singular transition approximations, without imposing continuity on the greedy actor. Residual accuracy and transition precision then enter the same explicit economic error budget, while the coupled resource program retains the cost of evaluating the stored successor.

The frozen shared-proposal experiment tests incremental query work on unseen states and an unseen dimension. It retains its conventional and polynomial comparators and its calibration-selected service. The ReLU proposal does not reduce query work on this catalogue. This direct test is distinct from the joint-domain construction theorem and from the independently evaluated complete-policy cost comparison. The article keeps all three objects separate rather than treating a valid certificate as proof of a computational or economic advantage.

\paragraph{Reading path.}
The article develops construction, accuracy, qualification, optimizer-class transfer and actual costs. Core proofs remain in the Technical Supplement. The Numerical and Development Appendix retains recursive-query, vector-control, reuse and historical evidence unchanged, followed by the new proofs and complete catalogues. Preserved companions retain diffusions, recursive and endogenous preferences, temporal selves and games under their own hypotheses; the expected-cost identities here are not silently extended to those objectives. Sections~\ref{sec:recursive86}--\ref{sec:complexity86} develop the complete-policy, population and horizon extensions. Section~\ref{sec:prospective86} connects them to prospective decisions without merging the observation units of the different experiments.

\section{The economic operator and its policy account}\label{sec:core49}
Consider dates $t=0,\ldots,T$ and an invariant compact state domain $X\subset\R^d$. At date $t$, the action belongs to a nonempty compact set $A_t(x)$. The transition is $F_t(x,a,Z_{t+1})$, the current cost is $c_t(x,a)$, and the terminal cost is $g(x)$. Write $B_0=1$ and $B_t=\prod_{j<t}\beta_j$, with $\beta_j>0$. An implemented policy includes its observation rule, action selection and feasible numerical realization, not just its stored parameters.

For a continuation $v$, define
\begin{equation}\label{eq:operator49}
 Q_t(v;x,a)=c_t(x,a)+\beta_t\rho_t(v(F_t(x,a,Z))),\qquad
 \mathcal T_tv(x)=\inf_{a\in A_t(x)}Q_t(v;x,a).
\end{equation}
Throughout the deterministic comparison, the evaluations $\rho_t$ are monotone and cash invariant on a common domain containing all functions and constant translates used below. They are therefore nonexpansive in the supremum norm. Continuous objectives, effective compact feasible correspondences and measurable minimizers, or measurable arbitrarily near-minimizers with their error charged, are assumed. The stipulated conditional recursion gives Bellman comparison against all admissible adapted controls. In the experiments $\rho_t$ is conditional expectation, costs are bounded and innovations have their original continuous uniform law. The direct-cost inference in Section~\ref{sec:decisions49} is specific to expectation; it is not automatically an inference theorem for nonlinear recursive utility or games.

A Neural Bellman Operator is a construction and evaluation procedure for \eqref{eq:operator49}. It forms targets using its \emph{own} fitted $f_{t+1}$, returns a neural continuation and a feasible policy, and records an economic error allowance. It does not receive the optimal continuation as a training label. Arbitrary trained networks can be evaluated by the first result below. The subsequent constructive backend specifies how a network and its actor are obtained without assuming successful nonconvex optimization.

\subsection{Action-dependent errors and complete-policy loss}
Let $H(a)$ be a true current cost and $\widehat H(a)$ its continuation approximation. If $\widehat a$ is $\delta$-optimal for $\widehat H$ on the same feasible action set, then
\begin{equation}\label{eq:center49}
 H(\widehat a)-\inf_aH(a)\le\delta+
 \operatorname{osc}_{a}\{\widehat H(a)-H(a)\}.
\end{equation}
Indeed, compare $\widehat a$ with an arbitrarily near-optimal true action and cancel the approximate objectives. The infimum limit gives the result. Additive continuation levels cancel; action-dependent discrepancies do not. Under a change in future policies, the continuation must be transported to the new economic problem before this inequality licenses reuse. The full development edition retains the original quantitative transport, recursive-domain and economic-application results. No transport bound is presumed valid after an unaccounted change of constraints or preferences.

\begin{proposition}[One-sided policy sandwich]\label{prop:sandwich49}
Suppose, uniformly in state,
\[
 f_t-\mathcal T_tf_{t+1}\le u_t,\qquad
 \mathcal T_t^\pi f_{t+1}-f_t\le v_t,
 \qquad \ell_T\le f_T-g\le u_T.
\]
Then the actual feasible policy satisfies
\begin{equation}\label{eq:sandwich49}
 0\le J_0^\pi-V_0^*\le
 \sum_{t<T}B_t(u_t+v_t)+B_T(u_T-\ell_T).
\end{equation}
The bound is invariant to date-specific constant shifts of the fitted continuations, provided the associated residual accounts are shifted consistently.
\end{proposition}
\noindent The proof is in Supplement, Section~\ref{supp:proof-core49-1}.


A sharper valid residual account for the same continuation and the same deployed actor sharpens knowledge of the policy. It does not retrain that policy or change its realized cost. Conversely, subtracting two upper bounds in \eqref{eq:sandwich49} does not identify the difference in actual costs. These distinctions govern both the historical comparisons and the new experiments.

\section{Constructing and compiling a feasible neural policy}\label{sec:construction49}
Assume the primitive objective has, for feasible state--action pairs, state modulus $A_t^Q=C_{x,t}+\beta_tF_{x,t}L_{t+1}$ and action modulus $D_t=C_{a,t}+\beta_tF_{a,t}L_{t+1}$. Let
\[
 d_H(A_t(x),A_t(z))\le K_t^A\|x-z\|_1,
 \qquad L_t=A_t^Q+D_tK_t^A,\qquad L_T=L_g.
\]
The Hausdorff condition and the primitive moduli imply that $\mathcal T_tf_{t+1}$ is $L_t$-Lipschitz whenever $f_{t+1}$ is $L_{t+1}$-Lipschitz. All moduli and numerical integration procedures are supplied by the economic primitives, not fitted from successful outcomes.

At state nodes $x_{t,i}$ with covering radius $h_t$, form effective feasible $k_t$-nets. Evaluate their own-future objectives with certified error at most $e_t$. Minimize the computed labels and retain the corresponding \emph{original} action $a_{t,i}$. A Borel repair $\Gamma_{t,i}(x)$ is feasible at $x$ and satisfies
\begin{equation}\label{eq:repair49}
 \|\Gamma_{t,i}(x)-a_{t,i}\|_1\le K_t^A\|x-x_{t,i}\|_1+\zeta_t.
\end{equation}
The allowance $\zeta_t$ concerns displacement; feasibility itself may not fail. Define
\begin{equation}\label{eq:cone49}
 f_t(x)=\min_i\{y_{t,i}+L_t\|x-x_{t,i}\|_1\},\quad
 i_t(x)=\min\arg\min_i\{y_{t,i}+L_t\|x-x_{t,i}\|_1\}.
\end{equation}
Return $\pi_t(x)=\Gamma_{t,i_t(x)}(x)$. Terminal labels approximate $g$ with error $e_T$ and use the same cone construction with modulus $L_g$.

\begin{theorem}[Own-future construction with feasible witnesses]\label{thm:construct49}
The functions in \eqref{eq:cone49} have exact affine--ReLU realizations and Lipschitz moduli $L_t$ independent of label error. The returned policy is feasible and measurable. It satisfies
\begin{align}\label{eq:enclosures49}
 -e_t\le f_t-\mathcal T_tf_{t+1}&\le D_tk_t+e_t+2L_th_t,\\
 \mathcal T_t^\pi f_{t+1}-f_t&\le e_t+D_t\zeta_t,\notag
\end{align}
and hence
\begin{equation}\label{eq:constructionloss49}
 J_0^\pi-V_0^*\le
 \sum_{t<T}B_t(D_tk_t+2e_t+2L_th_t+D_t\zeta_t)
 +B_T(2e_T+2L_gh_T).
\end{equation}
With effective covers, repairs and certified queries at the prescribed accuracy, allocating each nonnegative term a share of a positive economic tolerance gives a terminating finite construction. A fixed numerical precision or a capped finite experimental ladder need not contain that allocation.
\end{theorem}
\noindent The proof is in Supplement, Section~\ref{supp:proof-core49-2}.


The Lipschitz-extension principle is classical \citep{mcshane1934}. A native min-plus envelope and its exact ReLU circuit, using the same labels, witnesses and ties, are \emph{one method with two encodings}. They return the same function and policy. Neither Theorem~\ref{thm:construct49} nor a faster representation proves a uniquely neural capability or a training-success probability for Adam or L-BFGS. The mathematical increment is the complete feasible, own-future construction and implementation account; the empirical question is its work and policy value relative to conventional generators.

\subsection{Original-owner compilation on a tensor cover}
On a tensor cover $\mathcal X=\prod_{j=1}^d\{x_{j,0},\ldots,x_{j,n_j}\}$, put $S=\prod_j(n_j+1)$. For arbitrary labels $y_i$, including inconsistent inexact labels, compute at each grid node $v$
\[
 (z_v,o_v)=\min_{\rm lex,i}(y_i+L\|v-x_i\|_1,i).
\]
The index $o_v$ is the original label owner, not the transformed node $v$.

\begin{theorem}[Continuous original-owner identity]\label{thm:compile49}
For a closed tensor cell $C$ with corners $\mathcal V(C)$, and $x\in C$,
\begin{align}\label{eq:compile49}
 f(x)&=\min_{v\in\mathcal V(C)}(z_v+L\|x-v\|_1),\\
 (f(x),i_*(x))&=\min_{\rm lex,v\in\mathcal V(C)}
 (y_{o_v}+L\|x-x_{o_v}\|_1,o_v).\notag
\end{align}
Coordinatewise forward and backward relaxations compute $(z,o)$ in $O(dS)$ additions and comparisons and $O(S)$ storage. After locating the cell, a query uses $O(d2^d)$ arithmetic operations. Exact rational comparisons preserve all ties, including $L=0$. Using the original owner's repaired action preserves the uncompiled policy, its cost and its certificate exactly.
\end{theorem}
\noindent The proof is in Supplement, Section~\ref{supp:proof-core49-3}.


The nodal algorithm is the classical distance transform \citep{felzenszwalb2012}; preserving the original action through continuous off-grid queries is the policy obligation. If there are $J_{t,i}$ node actions and $M_t$ integration representatives, write $Q_t=M_t\sum_iJ_{t,i}$. Literal dense cone evaluation costs $O(\sum_t dS_{t+1}Q_t)$ in representation work. Compilation replaces this by $O(\sum_t[dS_{t+1}+Q_td2^d])$, plus cell location, primitive evaluation, integration, feasible repair, certificate and output. Rational bit lengths and state boxes crossing many cells incur additional costs. There is no dimension-free complexity assertion.

\subsection{Acquired states and conventional actors}
A $b$-bit coordinate sensor reports a dyadic cell and its midpoint $\widehat x$. On the unit cube, $\|x-\widehat x\|_1\le r(b)=d2^{-b-1}$. For the investment capacity $\kappa(x)$ below, clip the original action to the smallest capacity throughout that cell and round down to an action quantum $\Delta_t$. This is feasible at every state in the cell. Relative to an exact true-state repair, its displacement is at most $2K_t^Ar(b)+\Delta_t$. Transporting the witness cone from $\widehat x$ to $x$ adds at most $2L_tr(b)$, plus any certified selector objective error $\nu_t$. Thus the implemented-policy account adds
\begin{equation}\label{eq:acquire49}
 \sum_{t<T}B_t\{(2L_t+2D_tK_t^A)r(b)+D_t\Delta_t+\nu_t\}.
\end{equation}
No continuity of the discrete owner selector is assumed. At uncertain numerical comparisons the simulator retains all possible actors, rather than choosing a favorable one.

For conventional FVI the continuation is multilinear on its realized tensor grid. Its actual modulus is bounded by the largest coordinate edge slope. A nearest-node actor uses a feasible node minimizer and true-state repair. Let $h_t$ be the largest half-cell $\ell^1$ radius, $L_t^{\rm gr}$ the Bellman graph modulus, and $e_t$ the node query allowance. Its optimal residual upper bound is $D_tk_t+e_t+L_t^{\rm gr}h_t$. Its selected-policy upper bound is $e_t+E_t$, where
\[
 E_t=\sup_x\{y_{i(x)}+L_t^{\rm gr}\|x-x_{i(x)}\|_1-f_t(x)\}.
\]
Partitioning cells by nearest-node boundaries makes the expression separately affine in each coordinate. Its maximum is therefore at a partition vertex; the implementation encloses all such vertices. Both uniform and curvature-refined FVI receive this sharper one-sided account, not an artificially doubled baseline. Their acquired-state allowance is derived separately in \eqref{eq:fvisensor49}.

\section{Original-optimum accuracy in the investment economy}\label{sec:accuracy62}
The exact witness theorem answers an implementation question conditional on a continuation. An economist also needs to know how close the complete implemented policy is to the original optimum. We now establish that connection in the investment economy itself. The argument exploits a property not used by the earlier cell-range calculation: despite the bilinear transition, the Bellman objective remains jointly convex. This permits a second-order interpolation account and transfers verified nodal decisions to an implementable policy over the entire state domain. It does not assume that a fitted neural continuation is convex.

\paragraph{Relation to the preceding two-control query.}
The earlier query experiment used the fixed shared capacity $1/8$. The complete dynamic extension below uses the original state-dependent capacity $c(x)$, with the retained exposure columns, action-cost interaction and second shock. Its nesting claim concerns the original scalar economic model; it is not an assertion that a full dynamic policy was already present in the earlier isolated query.

\subsection{Primitives and Bellman regularity}
Let $d\ge2$ be even, let indices be cyclic, and set $\beta=15/16$. Write
\begin{align*}
 S(x)&=\frac2d\sum_j(x_j-5/8)^2+
       \frac1{4d}\sum_j(x_j-x_{j+1})^2+
       2\left(\frac12-\frac2d\sum_jx_j\right)_+^2,\\
 g(x)&=S(x)+\frac2d\sum_j(x_j-5/8)^2,\qquad
 c(x)=\frac18+\frac1{8d}\sum_jx_j.
\end{align*}
The original scalar control has $0\le a\le c(x)$, running cost $S(x)+a^2+4a^4$, and transition
\[
 F_j(x,a,z)=\frac1{16}+\frac{x_j}{2}+\frac{x_{j+1}}8+
       \frac{x_j(1-x_{j+1})}{16}+B_ja+s_jz,
\]
where $(B_j,s_j)=(1/2,1)$ at alternating coordinates and $(1/4,-1)$ at the others, and $z$ is uniform on $[-1/32,1/32]$.

For the retained two-control extension, $a_1,a_2\ge0$, $a_1+a_2\le c(x)$, the second exposure column interchanges $1/2$ and $1/4$, and the action cost is
\[
 R(a)=a_1^2+a_2^2+4(a_1^4+a_2^4)+a_1a_2/4.
\]
An independent common-direction uniform innovation of radius $1/64$ is added. This is the previously specified extension, nested at $a_2=0$ when that innovation is disabled. Both models preserve $[0,1]^d$. Let $m\in\{1,2\}$ denote the number of controls and define
\begin{align*}
 Q_t^*(x,a)&=S(x)+R(a)+\beta\mathbb E V_{t+1}^*(F(x,a,Z)),\\
 V_t^*(x)&=\min_{a\in A(x)}Q_t^*(x,a),\qquad V_T^*=g.
\end{align*}
All costs and laws in this display are the economic primitives, not a fitted surrogate.

\begin{theorem}[Regularity of the original Bellman problem]\label{thm:regularity62}
For either control specification and every finite horizon, $V_t^*$ is convex and continuously differentiable in the interior, with the corresponding one-sided extension on the cube. It is $G_t$-Lipschitz in the $\ell^1$ norm and $M_t$-semiconcave in the Euclidean norm, where valid constants are
\begin{align*}
 G_T&=10/d,&G_t&=\frac{1037}{128d}+\beta\frac{47}{64}G_{t+1},\\
 M_T&=26/d,&M_t&=\frac{22}{d}+\frac{21}{256d}
             +\beta\left(\frac9{16}M_{t+1}+\frac18G_{t+1}\right).
\end{align*}
In particular, $G_t\le27/d$ and $M_t\le54/d$. The function $Q_t^*$ is jointly strongly convex on its convex feasible graph, with state curvature at least $107/(128d)$ and action curvature at least $7/4$.

A valid semiconcavity constant along each coordinate line is
\begin{align*}
 D_T&=\frac9d+\frac{16}{d^2},\\
 D_t&=\frac5d+\frac{16}{d^2}+\frac{21}{256d^2}
 +\beta M_{t+1}\left(\frac{85}{256}+\frac{23}{256d}\right).
\end{align*}
Consequently, for the positive multilinear interpolant $I_n$ on a mesh of spacing $1/n$,
\begin{equation}\label{eq:interpolation62}
 0\le I_nV_t^*(x)-V_t^*(x)\le\kappa_t,
 \qquad \kappa_t=\frac{dD_t}{8n^2}.
\end{equation}
\end{theorem}

\noindent The proof is in Supplement, Section~\ref{supp:proof-accuracy62-print-1}.


The theorem is a structural statement about these economic primitives. It is not a claim that every controlled diffusion, recursive preference specification, or game in the complete development edition has a convex Bellman problem. Nor does it require a signed-weight neural fit to preserve convexity. Convex approximation of dynamic programs has an established literature; the additional work here is to verify the needed curvature for the original nonlinear capacity-constrained model and connect it to the actual NBO actor and complete error account.

\subsection{Certified nodal Bellman bounds}
Suppose nodal arrays satisfy $l_{t+1}\le V_{t+1}^*\le u_{t+1}$ at all grid vertices. Equation~\eqref{eq:interpolation62} gives the functions $I_nl_{t+1}-\kappa_{t+1}$ and $I_nu_{t+1}$ as valid lower and upper continuations. The lower function need not itself be convex; convexity is used for the unknown true value and for the comparison argument, not asserted of arbitrary endpoint arrays.

Divide each uniform innovation component into $q$ equal bins and use its conditional mean. Convexity gives a lower quadrature bound for the true continuation. Semiconcavity gives an upper error
\begin{equation}\label{eq:quadrature62}
 \zeta_{t+1}=\frac{M_{t+1}}{6q^2}
                      \sum_{r=1}^{s}\rho_r^2\|s_r\|_2^2,
\end{equation}
where $s=1$ or $2$, $\rho_1=1/32$, and $\rho_2=1/64$ in the extension. The expectation remains the original continuous law. At the terminal date $g$ is evaluated directly, so no terminal interpolation allowance is needed.

At a vertex $v$, let $\mathcal A_A(v)$ be the capacity simplex lattice with spacing $c(v)/A$, including all its faces. For $m=1,2$, a minimizer can be represented by neighboring lattice vertices whose coordinate variances sum to at most $m c(v)^2/(4A^2)$. A valid action smoothness constant is
\[
 H_t=\frac{21}{4}+\beta M_{t+1}\|B\|_2^2,
 \quad \|B\|_2^2=\frac{5d}{32}\ (m=1),\quad
 \|B\|_2^2=\frac{9d}{32}\ (m=2).
\]
Taylor's inequality and that barycentric representation imply
\begin{equation}\label{eq:action-cover62}
 0\le\min_{a\in\mathcal A_A(v)}Q_t^*(v,a)-V_t^*(v)
       \le\delta_t:=\frac{mH_t}{128A^2}.
\end{equation}
The linear Taylor terms cancel even when the optimum lies on a face. This is an error against the continuous feasible optimum, not only against the finite action menu.

Let $q_t^-(v,a)$ and $q_t^+(v,a)$ include all interpolation, quadrature and arithmetic errors and enclose $Q_t^*(v,a)$. Then
\begin{equation}\label{eq:nodal62}
 l_t(v)=\max\{0,\min_{a\in\mathcal A_A(v)}q_t^-(v,a)-\delta_t\},
 \qquad
 u_t(v)=\min_{a\in\mathcal A_A(v)}q_t^+(v,a)
\end{equation}
are valid nodal Bellman bounds. The zero truncation uses the nonnegative economic costs. Extra neural or conventional proposals may improve the upper side but do not remove the lower action-cover allowance. Positive interpolation, averaging and minimum maps are monotone and nonexpansive; the separate allowances in \eqref{eq:interpolation62}--\eqref{eq:action-cover62} are therefore propagated, rather than being discarded because an endpoint calculation contains the truth.

\subsection{From a fitted witness to the complete implemented policy}
Let $a_t(v)\in A(v)$ be the actual action returned at each vertex by any generator. It may be selected by the exact neural witness, a conventional fitted continuation, or their common verified fallback. Write $w_v(x)$ for the multilinear weights and set $\bar\pi_t(x)=\sum_vw_v(x)a_t(v)$. Affineness of capacity makes this action feasible. Let $q_t^{\mathrm{sel},+}(v)$ be a verified upper bound on its original-optimum continuation value $Q_t^*(v,a_t(v))$, using the same lower and upper Bellman reference for all generators.

\begin{theorem}[Actual-policy, whole-domain certificate]\label{thm:policy62}
Set
\[
 e_t=\max_v\{q_t^{\mathrm{sel},+}(v)-l_t(v)\}+\kappa_t.
\]
Suppose the implemented policy $\widehat\pi_t$ is feasible and differs from $\bar\pi_t$ by at most $b_t$ in action $\ell^1$ norm, uniformly over the state cube. Define
\[
 \Lambda_t=\frac{13}{16}+\beta\frac{3d}{8}G_{t+1},\qquad
 D_T^\pi=0,\quad D_t^\pi=e_t+\Lambda_tb_t+\beta D_{t+1}^\pi.
\]
Then, at every true state and date,
\[
 0\le J_t^{\widehat\pi}(x)-V_t^*(x)\le D_t^\pi.
\]
This statement is conditional on the actual recorded policy, not on success of its fitting algorithm. It applies equally to a pure fitted policy and to a policy augmented by verified common candidates.
\end{theorem}

\noindent The proof is in Supplement, Section~\ref{supp:proof-accuracy62-print-2}.


In the execution, observation is rounded down at forty bits. An acquired state $y\le x$ has no greater capacity than the true state. Each interpolated action is then rounded down at twenty bits using the lower endpoint of its arithmetic enclosure. If $L_{ij}$ bounds the derivative of action coordinate $i$ with respect to state coordinate $j$, and $\eta$ bounds native interpolation error, a valid implementation allowance is
\[
 b_t\le 2^{-40}\sum_{i,j}L_{ij}
           +m(2^{-20}+2\eta+\eta_{\mathrm{out}}).
\]
Here $L_{ij}$ is computed from the maximum adjacent nodal difference times $n$, and $\eta_{\mathrm{out}}$ covers endpoint subtraction. All these numbers are recorded outward. This new barycentric acquired actor is explicitly distinguished from the earlier discontinuous cell-constant actor; an old actor is not silently replaced during a purported fixed-policy revalidation.

\subsection{Arithmetic, work, and scope}
The native interpolation kernel forms nonnegative tensor weights and a weighted sum of at most $2^d$ stored binary64 endpoints. It is compiled without reassociation or floating-point contraction. For $K=(3d+4)2^d+32$, an explicit conservative absolute allowance is
\[
 \eta=\frac{K2^{-53}}{1-K2^{-53}}\max_v|z_v|+K2^{-1022}.
\]
The final term conservatively covers underflow. Input coordinates, affine indexing, finite endpoints and the binary64 format are checked. Roundoff expansion and the ordinary interval cost arithmetic remain separate from the analytic approximation bounds. Exact rational fixtures compare the native sum with its mathematical tensor interpolant.

With $N=(n+1)^d$ vertices, $K_A=A+1$ scalar actions or $K_A=(A+1)(A+2)/2$ two-control actions, the reference calculation uses $O(TNK_Aq^s2^d d)$ elementary work under the explicit tensor implementation. Stored reference values use $O(TN)$ words and a policy uses $O(TNm)$ words. Streaming batches avoid a state-pair table. Proposal fitting, exact rational witness work, actor transfer, serialization, independent cost inference and compilation are additional costs, not absorbed into this bound. The analytic errors are second order in state, action and innovation resolution; the state-count and tensor-evaluation factors still depend exponentially on dimension. The result supplies a matched original-optimum comparison without asserting that these costs have disappeared.

\section{Frozen neural continuations and policy-cost precision}\label{sec:frozen72}
The optimal-advantage estimator in Section~\ref{sec:centered69} makes a policy comparison informative, but its conditional optimal-value queries are expensive. A fitted continuation can replace those queries without changing the estimand. The necessary certificate concerns its Bellman residual across states, not its distance from an unknown optimizer. This places the original NBO continuation, rather than only a proposed action, inside the evaluation calculation.

\subsection{Residual oscillation is the relevant range}
Consider the bounded expected-cost model with common terminal cost $g$, discounts $\beta_t\ge0$, and a common initial-state law. Write $b_0=1$ and $b_t=\prod_{j<t}\beta_j$. Fix bounded measurable continuations $h_0,\ldots,h_T$, conditionally on their construction. Define
\begin{align*}
 Q_t^h(x,a)&=c_t(x,a)+\beta_tP_t^a h_{t+1}(x),\\
 r_t^h(x)&=\inf_{a\in A_t(x)}Q_t^h(x,a)-h_t(x),\\
 a_t^{h,\pi}(x)&=Q_t^h(x,\pi_t(x))-\inf_{a\in A_t(x)}Q_t^h(x,a).
\end{align*}
The infimum is over the original continuous feasible set. Suppose verified constants give $0\le a_t^{h,\pi}\le\eta_t^\pi$, and likewise for $\nu$, throughout the true state domain. Let $\omega_t$ bound the oscillation of $r_t^h$, and $\omega_T$ that of $g-h_T$. Set
\[
 \Omega_h=\sum_{t<T}b_t\omega_t+b_T\omega_T,
 \qquad H_h^\pi=\sum_{t<T}b_t\eta_t^\pi.
\]
These constants include acquisition, action rounding and arithmetic allowances wherever those operations affect the implemented policy or the verified functions. A small empirical residual alone is not such a certificate.

\begin{theorem}[Frozen-continuation comparison]\label{thm:frozen72}
Along the actual trajectory of $\pi$, put
\begin{equation}\label{eq:frozen-observation72}
 Z_h^\pi=\sum_{t<T}b_t\bigl[Q_t^h(X_t^\pi,\pi_t(X_t^\pi))-h_t(X_t^\pi)\bigr]
                 +b_T(g-h_T)(X_T^\pi).
\end{equation}
For any coupling of the two original trajectory laws, $D_h=Z_h^\pi-Z_h^\nu$ satisfies
\begin{align}
 \mathbb E D_h&=\mathbb E[J_0^\pi-J_0^\nu],\label{eq:frozen-identity72}\\
 -\Omega_h-H_h^\nu&\le D_h\le\Omega_h+H_h^\pi.\label{eq:frozen-range72}
\end{align}
Moreover, at every initial state,
\begin{equation}\label{eq:frozen-regret72}
 0\le J_0^\pi-V_0^*\le\Omega_h+H_h^\pi.
\end{equation}
Adding an arbitrary state-independent constant to each $h_t$ changes neither $D_h$ nor these oscillation bounds. With $h_t=V_t^*$, the result reduces to optimal-advantage centering.
\end{theorem}
\noindent The proof is in Supplement, Section~\ref{supp:frozen72}.

The terminal correction in \eqref{eq:frozen-observation72} is essential unless $h_T=g$. Conditional expectation, not pathwise equality with the realized cash-cost difference, supplies \eqref{eq:frozen-identity72}. The result uses the actual kernels and costs; it assumes neither a correct statistical transition model estimated from data nor a convex neural fit. The pointwise policy certificate also applies conditionally after an adaptive choice among valid continuations. The independent-sample confidence statement below requires its additional sampling conditions.

For $M$ conditionally fixed comparisons, suppose independent trajectory rows yield outward enclosures $[L_i,U_i]$ of $D_h$, intersected with its valid support. Write $R_h=2\Omega_h+H_h^\pi+H_h^\nu$. A simultaneous interval with failure probability at most $\alpha$ is
\begin{equation}\label{eq:frozen-ci72}
 \left[\overline L-R_h\sqrt{\frac{\log(2M/\alpha)}{2n}},\quad
       \overline U+R_h\sqrt{\frac{\log(2M/\alpha)}{2n}}\right].
\end{equation}
This follows by the bounded-sum inequality applied to both endpoint samples and a union bound. The numerical and sampling width allocations of Corollary~\ref{cor:precision69} apply with $R_h$. Fitting, uniform residual qualification, conditional integration and any subsequent selection are distinct paid operations. In particular, a cache that adapts across rows changes the controller's history: its rows cannot be treated as independent evaluations of one frozen Markov policy. A fixed cache or an independent episode-reset design is sufficient to retain the stated sampling construction.

\subsection{Conditional integration of an actual neural continuation}
A shallow rational affine--ReLU continuation has the form
\[
 h(y)=c+\ell^\top y+\sum_{j=1}^K v_j(w_j^\top y+u_j)_+.
\]
Its readout weights may have either sign. In the original investment laws the next state, conditional on $(x,a)$, is affine in one or two independent centered uniform innovations. A neuron therefore has projected center $z$ and nonnegative projection radii $A,B$. Write $z_+=\max\{0,z\}$.

\begin{proposition}[Exact conditional ReLU moments]\label{prop:moment72}
For independent uniform projections of radii $A,B>0$, the conditional means of a neuron are
\begin{align}
 \mathcal M_1(z;A)&=\frac{(z+A)_+^2-(z-A)_+^2}{4A},\label{eq:moment-one72}\\
 \mathcal M_2(z;A,B)&=\frac{1}{24AB}\bigl[(z+A+B)_+^3-(z+A-B)_+^3\nonumber\\
 &\hspace{37mm}-(z-A+B)_+^3+(z-A-B)_+^3\bigr].\label{eq:moment-two72}
\end{align}
A zero projection radius is removed; if all radii vanish the mean is $z_+$. Consequently $P^ah$ is a finite exact rational calculation for rational primitives, states, actions and weights. For $n$ trajectories and at most $K$ neurons per date, its arithmetic-operation count is $O(nTKd)$, exclusive of fitting, trajectory construction and uniform qualification. No optimal-value descendant query occurs in this calculation.
\end{proposition}
\noindent The proof is in Supplement, Section~\ref{supp:moments72}.

The operation count does not bound rational bit complexity. Near-zero projected radii can also make a naive floating evaluation unstable; directed rounding or rational evaluation, including the zero-radius branches, is required. The executed moment catalogue checks the rational identities against an independently written polygon integral. It is a conditional-integration experiment, not a new trained-policy cost comparison. The formulas do not imply that the achieved residual oscillation of an arbitrary fitted network is small.

\subsection{Qualifying a fitted continuation without a Bellman tree}
Uniform qualification need not query $V^*$. Suppose $Q_t^h$ is Lipschitz in state and action with constants $L_{x,t},L_{a,t}$, and $h_t$ has constant $L_{h,t}$, all in $\ell^1$. In the investment capacity simplex, an action feasible at $y$ can be repaired continuously at $x$ with displacement at most $\kappa\|x-y\|_1$, where $\kappa=1/(8d)$. Thus a valid residual Lipschitz constant is
\[
 L_{r,t}=L_{x,t}+\kappa L_{a,t}+L_{h,t}.
\]
Let a finite state cover have fill radius $\rho$, and let $[l_j,u_j]$ enclose $r_t^h$ at every cover point. Then
\begin{equation}\label{eq:residual-cover72}
 \min_jl_j-L_{r,t}\rho\le r_t^h(x)
       \le\max_ju_j+L_{r,t}\rho
\end{equation}
throughout the cube. At a cover point, a feasible action mesh of fill radius $s$ gives a lower infimum bound by subtracting $L_{a,t}s$ from the smallest action-value lower endpoint; the smallest upper endpoint is an upper bound. Exact moments can supply those action endpoints. None of these steps requires convexity of $h$ or a true-value curvature minorant for $Q^h$.

For $N_t$ state-cover points and $A_t$ actions per point, this qualification costs $O(\sum_tN_tA_tK_{t+1}d)$ arithmetic operations, plus certification of the moduli and any actor-specific greedy-gap bound. Its cover can still be expensive in high dimension. The result replaces a recursive optimal-value oracle by an explicit approximation-and-qualification tradeoff; it does not make that qualification free. The residual-cover proof, the action-mesh argument and the terminal allowance are given in the supplement. Value-function control variates have an established variance-reduction literature \citep{munos2006variance}; the additional requirement here is an original-law, whole-domain certificate for the fitted NBO continuation and the actual acquired policy.

\section{Fitting a continuation for qualification}\label{sec:qualification73}
A continuation used to evaluate a policy should be fitted to the error that determines the evaluation cost. The residual-cover bound in \eqref{eq:residual-cover72} supplies such an objective. We now make its optimization constructive. The future continuation is held fixed at each backward step; neither an optimal-value oracle nor success of unrestricted hidden-layer training is assumed.

Fix a date and suppress its subscript. Let $x_1,\ldots,x_N$ be a state cover of radius $\rho$ in the $\ell^1$ metric, and suppose
\[
 y_j^-\le (\mathcal T h_+)(x_j)\le y_j^+.
\]
These are verified bounds on the original continuous-action infimum, not only on a proposal menu. Write $h_\theta=\sum_{k=1}^K\theta_k\phi_k$, where the fixed features have verified Lipschitz constants $L_k$ and $|\theta_k|\le B_k$. The constant feature is allowed with $L_k=0$. Let $C$ be a verified Lipschitz constant of $\mathcal T h_+$, including state-dependent feasibility repair. With $F_{jk}=\phi_k(x_j)$, define
\begin{equation}\label{eq:readout-objective73}
 \mathcal W(\theta)=\max_j(y_j^+-F_j\theta)
       -\min_j(y_j^--F_j\theta)
       +2\rho\left(C+\sum_k L_k|\theta_k|\right).
\end{equation}

\begin{theorem}[Qualification-optimal readout]\label{thm:readout73}
For every feasible $\theta$, $\operatorname{osc}(\mathcal T h_+-h_\theta)\le\mathcal W(\theta)$. Minimizing \eqref{eq:readout-objective73} is a finite linear program with $2K+2$ variables and $2N+2K$ inequality constraints, apart from variable bounds. For probability vectors $p,q$ on the nodes, set $\lambda_k=2\rho L_k$ and
\begin{equation}\label{eq:readout-dual73}
 \mathcal D(p,q)=p^\top y^+-q^\top y^-+2\rho C
   -\sum_k B_k\bigl(|[F^\top(q-p)]_k|-\lambda_k\bigr)_+.
\end{equation}
Then $\mathcal D(p,q)\le\inf_\theta\mathcal W(\theta)$, with equality for some $p,q$. Thus any feasible readout and exact probability vectors certify an achieved optimization gap of at most $\mathcal W(\theta)-\mathcal D(p,q)$.
\end{theorem}

The proof is in the Technical Supplement. It constructs the primal constraints and dual support function explicitly. Rationalizing a solver's proposed readout and probability vectors gives a checkable certificate; a convergence flag is not the certificate. Label intervals, features, moduli, coefficient bounds and the cover remain part of the numerical contract.

Applying the theorem backwards gives a fitted continuation at every date. Its achieved ranges $\mathcal W_t$, together with the terminal oscillation and the actual implemented greedy-gap bounds, can be substituted directly into Theorem~\ref{thm:frozen72}. This connects the attained fitting objective to both original-optimum regret and the range of the original-law cost observation. Each readout minimizes its stated one-step qualification bound conditional on the fixed future. It does not minimize the joint multi-date range, the true policy loss, or a hidden-feature objective. The next backward step recomputes its labels and moduli from the continuation actually returned.

The label construction in \eqref{eq:residual-cover72} uses exact conditional moments for a shallow ReLU future and a verified continuous-action cover. Its work is charged in addition to the linear program, coefficient rationalization, modulus calculation and dual check. A conventional basis with the same number of coefficients receives the same labels and the same objective; it uses its own valid feature moduli. This is a comparison of qualified approximations, not an assertion that different bases span the same functions or that a neural encoding has an intrinsic advantage.

\section{Representation calibration and optimization-stable policies}\label{sec:calibration81}
The attained continuation in Section~\ref{sec:frozen72} is an economic object, not merely an optional action proposal. Its qualification must charge the represented function rather than an avoidable featurewise overestimate. We make this comparison explicit and then pass from an attained fit to a class of optimization outputs with the same implemented policy. All statements in this section retain the original two-date investment law, fixed terminal payoff, acquired cells and feasible action menus.

\subsection{A common affine calibration}
Write a two-state additive continuation as
\[
 h_\theta(x)=\theta_0+\sum_{j=1}^4\theta_{1j}\psi_j(x_1)
                         +\sum_{j=1}^4\theta_{2j}\psi_j(x_2).
\]
For the ReLU basis, $\psi_j(s)=(s-(j-1)/4)_+$. For the Bernstein basis, $\psi_j(s)=\binom4j s^j(1-s)^{4-j}$; the omitted zeroth coefficient on each coordinate is zero. Both representations have nine coefficients and use the same coefficient box. Define
\begin{align}
 L_R(\theta)&=\max_{i=1,2}\max_{k=1,\ldots,4}
                      \left|\sum_{j=1}^k\theta_{ij}\right|,\label{eq:relu-modulus81}\\
 L_B(\theta)&=4\max_{i=1,2}\max_{j=1,\ldots,4}
                      |\theta_{ij}-\theta_{i,j-1}|,
                      \qquad \theta_{i0}=0.\label{eq:bernstein-modulus81}
\end{align}
These are state Lipschitz moduli for the $\ell^1$ distance. The first is exact for the represented piecewise-affine function; the second is a derivative-control bound and need not be exact for a general polynomial.

Suppose the fixed-label Bellman value at nodes $x_i$ belongs to $[y_i^-,y_i^+]$, every state is within $\ell^1$ distance $\rho$ of a node, and $G$ bounds the Bellman target's state modulus. With $\phi_i^\top\theta=h_\theta(x_i)$, put
\begin{equation}\label{eq:qualification81}
 F(\theta)=\max_i(y_i^+-\phi_i^\top\theta)
       -\min_i(y_i^--\phi_i^\top\theta)+2\rho\{G+L(\theta)\}.
\end{equation}
Here the labels, $G$, the cover and the feature map are fixed before this readout is optimized. A change in the fitted future defines a new fixed-label problem.

\begin{proposition}[Affine-exact qualification]\label{prop:calibration81}
The moduli in \eqref{eq:relu-modulus81}--\eqref{eq:bernstein-modulus81} give valid residual enclosures and the residual-oscillation upper bound \eqref{eq:qualification81}. The objective is convex and admits a finite linear-programming representation. For every common affine continuation $h(x)=a+u x_1+v x_2$, whenever its two coefficient representations are feasible, both moduli equal $\max\{|u|,|v|\}$. Thus, for common labels, node cover and target modulus, the two qualification objectives agree on this affine function.
\end{proposition}

\begin{proof}
A ReLU coordinate has the displayed cumulative slopes on its four affine intervals. The dual norm of $\ell^1$ is $\ell^\infty$, so their largest coordinate magnitude gives the state modulus. A Bernstein derivative is a convex combination of $4(\theta_{ij}-\theta_{i,j-1})$ with degree-three Bernstein weights. This proves its bound. Residuals change by at most $(G+L)\rho$ between a node and a covered state, proving the enclosure and its oscillation bound. Each term in \eqref{eq:qualification81} is a maximum of finitely many affine functions. For an affine function, the ReLU coefficients are $(u,0,0,0)$ and $(v,0,0,0)$; the Bernstein coefficients are $(u/4,u/2,3u/4,u)$ and the analogous vector for $v$. All relevant slopes or derivative controls are $u$ or $v$.
\end{proof}

This calibration does not identify the two non-affine function classes, equalize every numerical conditioning effect, or remove all conservatism from Bernstein controls. It isolates a particular qualification difference. For example, the old coefficientwise modulus charges $ux_1$ as $|u|$ in its ReLU encoding but as $10|u|$ in its Bernstein encoding. The calibrated modulus charges both as $|u|$. The factorial experiment below retains the old fits and evaluates the new fits separately; the calibration is not an ex post correction of their recorded costs.

\subsection{Exact certificates for action stability}
Fix the terminal-label qualification problem and an attained coefficient vector $\widehat\theta$, with $U=F(\widehat\theta)$. On each first-date observation cell $C$, let $a_C$ be its attained robust greedy action, evaluated at the declared cell representative. The terminal policy is fixed by the common terminal payoff. For each different action $b$ in the same robust menu, exact conditional integration gives an affine contrast
\[
 g_{Cb}(\theta)=Q^{h_\theta}(x_C,b)-Q^{h_\theta}(x_C,a_C)
                       =c_{Cb}+d_{Cb}^\top\theta.
\]
The switch problem minimizes $F(\theta)$ on the coefficient box subject to $g_{Cb}(\theta)\le0$. Let $D_{Cb}$ be any verified lower bound on its infimum. If the constraint is infeasible, its infimum is $+\infty$. Define $D_C=\min_{b\ne a_C}D_{Cb}$ and
\begin{equation}\label{eq:stable-cells81}
 \mathcal S_\epsilon=\{C:D_C>U+\epsilon\},\qquad \epsilon\ge0.
\end{equation}
Strict inequality excludes ties as well as strict reversals.

\begin{theorem}[Optimization-to-policy transfer]\label{thm:stability81}
Every feasible readout satisfying $F(\theta)\le U+\epsilon$ selects exactly $a_C$ on every cell in $\mathcal S_\epsilon$. In particular this holds for every $\epsilon$-optimal solution of the fixed-label qualification problem, because its optimum is at most $U$.

Let $q_\epsilon$ be the initial-law probability of the cells outside $\mathcal S_\epsilon$. Suppose the total cash cost of every such complete feasible controller belongs to $[0,H]$. With the terminal policy fixed as above,
\begin{equation}\label{eq:cost-stability81}
 \sup_{F(\theta)\le U+\epsilon}
       |J^{\pi_\theta}-J^{\pi_{\widehat\theta}}|
                       \le Hq_\epsilon.
\end{equation}
If all cells are certified, every readout in that sublevel gives the same complete policy and the same original-law cost. Under the uniform initial law on $[0,1]^2$ and an $n$-by-$n$ equal-cell cover, $q_\epsilon$ is exactly the fraction of uncertified cells; boundary assignments do not change that probability.
\end{theorem}

\begin{proof}
If $b$ weakly displaced $a_C$, the readout would be feasible for its switch problem and hence $F(\theta)\ge D_{Cb}\ge D_C$. This contradicts $F(\theta)\le U+\epsilon<D_C$. Couple the attained and alternative controllers with the same initial state and innovations. On a certified initial cell their first actions agree, their next states agree, and their fixed terminal-policy actions agree. Their cash costs are identical there. Elsewhere the absolute cash-cost difference is at most $H$, giving \eqref{eq:cost-stability81} by integration.
\end{proof}

The theorem does not replace failed switch certificates by successful ones: a nonpositive surplus is unresolved. Nor does a terminal-label certificate automatically propagate through newly fitted label problems at a longer horizon. The numerical result here concerns exactly the two-date construction for which only $h_1$ determines the first-date action and the terminal actor is common.

The switch bounds have an exact dual representation. Write $L(\theta)=\max_j|s_j^\top\theta|$, coefficient limits $|\theta_k|\le B_k$, and $\alpha=2\rho$. For probability vectors $p,q$, nonnegative multipliers $\xi_j^+,\xi_j^-$ with $\sum_j(\xi_j^++\xi_j^-)\le\alpha$, and $\nu\ge0$, set
\begin{align}
 z&=\sum_i(q_i-p_i)\phi_i
       +\sum_j(\xi_j^+-\xi_j^-)s_j+\nu d_{Cb},\nonumber\\
 D(p,q,\xi,\nu)&=p^\top y^+-q^\top y^-+2\rho G
                        +\nu c_{Cb}-\sum_kB_k|z_k|.\label{eq:switch-dual81}
\end{align}
Every such number is a valid $D_{Cb}$. The proof and independent exact replay are given in Appendix~\ref{app:dual81}. A floating optimizer supplies proposed multipliers only; their nonnegativity, normalization, total mass and lower-bound value are verified rationally.

\subsection{Simultaneous economic conclusions for optimizer classes}
Suppose the frozen attained-policy comparison supplies a simultaneous interval $[l_{jk},u_{jk}]$ for $J^{\pi_j}-J^{\pi_k}$. Apply Theorem~\ref{thm:stability81} separately to their fixed-label sublevels, with radii $v_j=H_jq_j$ and $v_k=H_kq_k$. Then, on the same event,
\begin{equation}\label{eq:uniform-comparison81}
 J^{\pi_\theta}-J^{\pi_\vartheta}
       \in[l_{jk}-v_j-v_k,\ u_{jk}+v_j+v_k]
\end{equation}
simultaneously for every pair of readouts in those sublevels. No union bound over an uncountable optimizer class is needed: the extension is deterministic on the original finite-policy event. When both covers are entirely stable the interval is unchanged. Sharp controller-choice and deployment-allocation results in Section~\ref{sec:allocation79} can therefore be applied directly to these qualified optimizer classes after the stated enlargement. This transfers attained-cost evidence without identifying numerical fitting success with either population training reliability or a new economic observation.

\section{Economic accuracy of approximately optimized continuations}\label{sec:actionwise82}
A qualification tolerance need not preserve every acquired action to preserve an economic conclusion. The bound in Theorem~\ref{thm:stability81} charges the full cost range on each uncertified cell. We now retain the particular actions that remain possible and the direction of their economic effects. This gives a cost guarantee for an entire optimization sublevel, including approximately optimal readouts, without treating an inconclusive switch certificate as an observed policy failure.

Fix a finite first-date acquisition partition $\mathcal C$, its initial-law masses $p_C$, a robust finite action menu on each cell, and a common feasible continuation policy $\tau$ after that date. These objects are held fixed throughout the comparison. The horizon may be any finite horizon for which the conditional costs below are enclosed. Use the fixed-label qualification objective $F$, attained readout $\widehat\theta$, value $U=F(\widehat\theta)$, reference action $a_C$, and switch lower bounds $D_{Cb}$ from Section~\ref{sec:calibration81}. For $\epsilon\ge0$, define
\begin{equation}\label{eq:possible-actions82}
 \mathcal A_C(\epsilon)=\{a_C\}\cup
       \{b\ne a_C:D_{Cb}\le U+\epsilon\}.
\end{equation}
The weak inequality deliberately retains ties. Let $\pi_\theta$ use a minimizing critic action at each cell representative and then $\tau$. Write $Q^\tau(x,b)=c_0(x,b)+\beta_0P_0^bJ_1^\tau(x)$, and, on every positive-mass cell, define the original-law contrast
\[
 \Delta_C(b)=\mathbb E[Q^\tau(X_0,b)-Q^\tau(X_0,a_C)\mid X_0\in C].
\]
Suppose a finite calculation provides $\ell_{Cb}\le\Delta_C(b)\le u_{Cb}$ for every retained action, with both endpoints zero for $a_C$. Set
\begin{equation}\label{eq:class-band82}
 \underline v_\epsilon=\sum_Cp_C\min_{b\in\mathcal A_C(\epsilon)}\ell_{Cb},
 \qquad
 \overline v_\epsilon=\sum_Cp_C\max_{b\in\mathcal A_C(\epsilon)}u_{Cb}.
\end{equation}
Zero-mass cells make no contribution.

\begin{theorem}[Actionwise optimization-to-cost transfer]\label{thm:actionwise82}
Under the preceding conditions, every feasible $\theta$ satisfying $F(\theta)\le U+\epsilon$ obeys
\begin{equation}\label{eq:uniform-cost82}
 \underline v_\epsilon\le J^{\pi_\theta}-J^{\pi_{\widehat\theta}}
                         \le\overline v_\epsilon.
\end{equation}
In particular the conclusion holds for every $\epsilon$-optimal readout. If all complete costs are in $[0,H]$ and the actionwise intervals are intersected with $[-H,H]$, then
\[
 -Hq_\epsilon\le\underline v_\epsilon\le0\le
       \overline v_\epsilon\le Hq_\epsilon,
\]
where $q_\epsilon$ is the initial mass of cells retaining a nonreference action.
For two such optimizer classes, an attained-policy simultaneous interval $[l_{jk},u_{jk}]$ gives, on the same event, the classwide interval
\begin{equation}\label{eq:class-pair82}
 [\,l_{jk}+\underline v_j-\overline v_k,\quad
       u_{jk}+\overline v_j-\underline v_k\,].
\end{equation}
These conclusions require no probability bound over the readouts in the sublevels.
\end{theorem}

\begin{proof}
Any critic-minimizing action $b\ne a_C$ satisfies $g_{Cb}(\theta)\le0$, so $D_{Cb}\le F(\theta)\le U+\epsilon$. It therefore belongs to \eqref{eq:possible-actions82}. Conditioning on the initial cell expresses the policy-cost difference as $\sum_Cp_C\Delta_C(b_C(\theta))$, which is between the endpoints in \eqref{eq:class-band82}. The infimum of $F$ is at most $U$, giving the optimization conclusion. A cell retaining only $a_C$ contributes zero; the signed support bounds every other contribution. Finally add and subtract the two attained costs to obtain \eqref{eq:class-pair82}. All steps after the attained-policy event are deterministic.
\end{proof}

The cellwise extrema relax the common-coefficient restriction and need not be attained by one readout. The theorem is consequently not a sharp joint characterization of the optimizer class. Its useful feature is different: it can retain an identified economic comparison when some actions remain uncertified and the binary $Hq_\epsilon$ allowance is too large. The original Bellman optimum and the constructive accuracy results are unchanged. A reference expected-loss upper bound can also be enlarged by $\overline v_\epsilon$; the actionwise comparison itself is not a new bound obtained by replacing the original optimum with a menu optimum.

For the executed two-date economy, the conditional contrasts are computed from the original stage costs, transitions and fixed terminal actor. Common innovations preserve their actual marginal laws. Centered polynomial differences retain cancellation, and rectangle queries enclose every crossed acquired cell. The final innovation is integrated analytically. Appendix~\ref{app:actionwise82} proves whole-bin validity, deterministic refinement convergence under vanishing arithmetic error, and the work bound. It also separates integration accuracy from the independent statistical interval to which \eqref{eq:class-pair82} is applied.

\section{Economically sufficient optimization accuracy}\label{sec:frontier83}
The qualification objective and the economic loss have different units. A numerical stopping rule should therefore be connected to the decision it supports, rather than chosen by equating their tolerances. The actionwise transfer supplies this connection for every tolerance in an interval, not only for a finite list of optimizer outputs.

For construction $j$, fix the labels, attained value $U_j$, action-switch bounds and continuation policy of Theorem~\ref{thm:actionwise82}. Fix $e_j>0$ and enclose every action with $D_{jCb}\le U_j+e_j$. Write $[v_j^-(\epsilon),v_j^+(\epsilon)]$ for its class-cost band. The only possible changes of this band on $[0,e_j]$ occur at
\[
 \mathcal K_j=\{0,e_j\}\cup
       \{D_{jCb}-U_j:0<D_{jCb}-U_j\le e_j\}.
\]
At a knot the weak switch inequality is already active. Thus the band is constant on each interval between successive knots, with the left endpoint included and the right endpoint excluded; the final endpoint is evaluated separately if it is itself a switch knot. No interpolation across a knot is permitted.

Let $[l_{jk},u_{jk}]$ be an attained-cost interval on the original simultaneous event. Define
\begin{align}
 L_{jk}(\epsilon,\eta)&=l_{jk}+v_j^-(\epsilon)-v_k^+(\eta),\label{eq:lower-frontier83}\\
 U_{jk}(\epsilon,\eta)&=u_{jk}+v_j^+(\epsilon)-v_k^-(\eta).\label{eq:upper-frontier83}
\end{align}
For an additional charge $\tau\ge0$ for construction $k$, and a declared economic equivalence margin $m>0$, put
\begin{align*}
 \mathcal R_{jk}(\tau)&=\{(\epsilon,\eta):L_{jk}(\epsilon,\eta)>\tau\},\\
 \mathcal E_{jk}(m)&=\{(\epsilon,\eta):-m<L_{jk}(\epsilon,\eta),\ U_{jk}(\epsilon,\eta)<m\}.
\end{align*}
Both sets are restricted to $[0,e_j]\times[0,e_k]$. A common-tolerance comparison is one diagonal slice; different constructions need not receive identical tolerances.

\begin{theorem}[Economic stopping frontier]\label{thm:frontier83}
On the attained-policy simultaneous event, every pair of feasible readouts satisfying
\[
 F_j(\theta)\le U_j+\epsilon,\qquad F_k(\vartheta)\le U_k+\eta
\]
has cost difference in $[L_{jk}(\epsilon,\eta),U_{jk}(\epsilon,\eta)]$, simultaneously for all tolerances in the stated rectangle. Consequently $\mathcal R_{jk}(\tau)$ certifies a strictly beneficial replacement by $k$ after its additional charge, and $\mathcal E_{jk}(m)$ certifies a cost difference strictly within the margin for the entire pair of readout classes.

These two sufficient regions are downward closed in their tolerance coordinates and are finite unions of knot rectangles with the stated endpoint conventions. Their maximal extent relative to the supplied actionwise bands can be obtained without a further optimizer call or cost observation. In particular, if a feasible numerical output has a verified primal--dual gap at most $\epsilon$, it belongs to the corresponding class and may be returned whenever its gap pair is in the required region. Choosing the stopping tolerance after examining the simultaneous intervals does not require an additional probability allowance.
\end{theorem}

\begin{proof}
As a tolerance increases, its possible-action set expands. Hence $v_j^-$ is nonincreasing and $v_j^+$ is nondecreasing. Each set changes only at the displayed switch knots, where equality retains the newly possible actions. Theorem~\ref{thm:actionwise82} gives \eqref{eq:lower-frontier83}--\eqref{eq:upper-frontier83} for every readout and every tolerance by the same deterministic argument. The lower endpoint decreases and the upper endpoint increases in either tolerance. This proves downward closure and the finite-rectangle representation, including the strict economic tests at their boundaries.

If $d_j\le\inf F_j$ is a verified dual value and $F_j(\theta)-d_j\le\epsilon$, then $F_j(\theta)\le\inf F_j+\epsilon\le U_j+\epsilon$. The same argument applies to $k$. All statements therefore hold for any stopping output satisfying these verified inequalities on the original event. Nothing in this implication requires a union bound over the outputs or tolerances.
\end{proof}

A boundary of a sufficient region is not an identified failure boundary of the true optimizer class. The cellwise relaxation can combine actions that no common coefficient vector selects. Nor does a strict region necessarily include its supremum tolerance: at an exact switch knot, the added action can destroy the desired strict economic inequality. Appendix~\ref{app:frontier83} gives an exact finite algorithm, checks both sides of every knot, and separates this stopping guarantee from a distributional claim about neural fitting success.

\section{Common readouts and economic policy qualification}\label{sec:compatible84}
An approximately optimized continuation is one function. Its coefficients determine the action in every observation cell. Separate action-switch tests deliberately forget part of this structure: an alternative may defeat the reference action without winning the complete feasible menu, and individually possible choices need not be jointly generated by the same readout. We retain these restrictions when transferring the original NBO qualification account to an economic decision. The attained continuation, economic primitives, action menu and policy-cost observations are unchanged.

\subsection{Winning actions and compatible choices}
Fix the finite acquisition partition $\mathcal C$, its original-law masses $p_C$, and a robust finite action menu $\mathcal A_C$. Each action in $\mathcal A_C$ is feasible throughout $C$. Let $\Theta=\prod_{j=1}^d[-B_j,B_j]$ be the declared readout box. With fitted labels and continuation policy $\tau$ held fixed, the implemented first-date rule minimizes
\[
 q_{Cb}(\theta)=c_{Cb}+v_{Cb}^{\mathsf T}\theta,
 \qquad b\in\mathcal A_C,
\]
using the original deterministic tie order. This is the fixed-feature own-future readout used in the preceding qualification construction, not an additional policy-value oracle. Write $F$ for its qualification objective, $\widehat\theta$ for the attained feasible readout, $U=F(\widehat\theta)$, and
\[
 \Theta_\epsilon=\{\theta\in\Theta:F(\theta)\le U+\epsilon\},
 \qquad \epsilon\ge0.
\]
The reference action in $C$ is $a_C$. The closed winner region
\begin{equation}\label{eq:winner84}
 W_{Cb}=\{\theta\in\Theta:q_{Cb}(\theta)\le q_{Ca}(\theta)
                     \text{ for every }a\in\mathcal A_C\}
\end{equation}
contains every readout that actually selects $b$. It can additionally contain ties resolved in favor of another action. Using this closed enlargement preserves the original tie rule and is conservative at its boundaries.

Suppose $D_{Cb}$ is a verified lower bound on $\inf_{W_{Cb}}F$. A separately certified empty winner region has bound $+\infty$. Retain the reference action and every alternative with $D_{Cb}\le U+\epsilon$ in $\mathcal A_C(\epsilon)$. An inherited reference-only switch bound remains valid and may be intersected with the full-menu bound by taking their maximum.

A joint event $H$ specifies actions $(C,b)$ at distinct cells. Suppose
\[
 D_H\le\inf\{F(\theta):\theta\in\bigcap_{(C,b)\in H}W_{Cb}\}
\]
is independently certified. If $D_H>U+\epsilon$, no member of $\Theta_\epsilon$ can select all the actions in $H$. Equality does not exclude the event. The collection of such verified events is denoted by $\mathcal H_\epsilon$; it need not contain every incompatible event.

Let $Q^\tau(x,b)$ be the actual original-law cost of taking $b$ and then following the common fixed continuation $\tau$. For each positive-mass cell, suppose numerical integration gives
\begin{equation}\label{eq:conditional-cost84}
 \ell_{Cb}\le
 \mathbb E[Q^\tau(X,b)-Q^\tau(X,a_C)\mid X\in C]
 \le u_{Cb}.
\end{equation}
Set $\ell_{Ca_C}=u_{Ca_C}=0$. Zero-mass cells contribute zero. All acquisition ambiguity, continuous-law integration and arithmetic allowances belong to these endpoints. The preceding actionwise calculation provides precisely these inputs.

\subsection{A computable economic bound}
For $z=(z_{Cb})$, define the assignment relaxation
\begin{equation}\label{eq:assignment84}
 \begin{split}
 \mathcal Z_\epsilon=\bigg\{z\ge0:\;&
 \sum_{b\in\mathcal A_C(\epsilon)}z_{Cb}=1\quad(C\in\mathcal C),\\
 &\sum_{(C,b)\in H}z_{Cb}\le |H|-1
                       \quad(H\in\mathcal H_\epsilon)\bigg\}.
 \end{split}
\end{equation}
Absent variables have value zero. Every policy induced by one readout in $\Theta_\epsilon$ defines a feasible zero--one assignment. Fractional assignments, and some unattainable zero--one assignments, can remain in $\mathcal Z_\epsilon$.

For any weights $w_{Cb}$, take numbers $\alpha_C\in\mathbb R$ and $\lambda_H\ge0$ satisfying
\begin{equation}\label{eq:assignment-dual84}
 \alpha_C+\sum_{H\ni(C,b)}\lambda_H\ge w_{Cb}
              \quad\bigl(b\in\mathcal A_C(\epsilon)\bigr),
\end{equation}
and define
\[
 V(w;\alpha,\lambda)=\sum_C\alpha_C+
                      \sum_{H\in\mathcal H_\epsilon}(|H|-1)\lambda_H.
\]
The dual need not be numerically optimal. Its feasibility and value can be checked exactly.

\begin{theorem}[Compatible optimizer-class cost bound]\label{thm:compatible84}
Under the preceding feasibility, winner, joint-event and conditional-cost conditions, let $V_\epsilon^+$ be any feasible dual value for $w_{Cb}=p_Cu_{Cb}$, and let $V_\epsilon^-$ be any feasible dual value for $w_{Cb}=-p_C\ell_{Cb}$. Then every $\theta\in\Theta_\epsilon$ satisfies
\begin{equation}\label{eq:compatible-cost84}
 -V_\epsilon^-\le J(\pi^\theta)-J(\pi^{\widehat\theta})
                    \le V_\epsilon^+.
\end{equation}
The exact assignment optima cannot give a wider interval than the full-menu cellwise extrema, which cannot give a wider interval than the inherited reference-only actionwise bound. Intersecting independently valid computed bounds preserves this ordering even when numerical dual proposals are not optimal.

For two constructed controllers $A,B$, suppose the attained-cost event gives
$L\le J(\pi^{\widehat\theta_A})-J(\pi^{\widehat\theta_B})\le U_0$.
Then all members of their respective qualified classes satisfy
\begin{equation}\label{eq:compatible-comparison84}
 L-V_A^- -V_B^+\le J(\pi^{\theta_A})-J(\pi^{\theta_B})
                       \le U_0+V_A^+ +V_B^-.
\end{equation}
The implication holds on that same event simultaneously for all allowances and all valid certificates, including certificates chosen after inspecting the covered intervals. It adds no failure probability for the continuum of readouts. A feasible numerical output with verified dual value $d\le\inf_\Theta F$ and $F(\theta)-d\le\epsilon$ belongs to $\Theta_\epsilon$.
\end{theorem}

\begin{proof}
Fix $\theta\in\Theta_\epsilon$. Its chosen action at each cell lies in the closed winner region. A singleton bound strictly larger than $U+\epsilon$ therefore cannot exclude that action. Similarly, selecting all members of a joint event in $\mathcal H_\epsilon$ would imply $F(\theta)\ge D_H>U+\epsilon$, a contradiction. Its choice indicators $z^\theta$ consequently belong to \eqref{eq:assignment84}.

The common continuation and the original initial law give the exact decomposition
\[
 J(\pi^\theta)-J(\pi^{\widehat\theta})
 =\sum_Cp_C\sum_bz^\theta_{Cb}
   \mathbb E[Q^\tau(X,b)-Q^\tau(X,a_C)\mid X\in C].
\]
Apply the upper inequality in \eqref{eq:conditional-cost84}. Multiplying \eqref{eq:assignment-dual84} by $z^\theta_{Cb}$, summing, and applying the nonnegative joint-event multipliers bounds the resulting expression by $V_\epsilon^+$. Apply the same argument to the negative lower endpoints for the other inequality in \eqref{eq:compatible-cost84}.

Removing all joint-event restrictions gives independent cell simplices. Their support is the sum of the cellwise maximal weights. Adding valid restrictions can only reduce that support. Full-menu exclusions retain a subset of the alternatives admitted by a reference-only comparison, proving the second ordering. Intersections preserve containment because each input bound contains the same cost difference.

Subtract the two instances of \eqref{eq:compatible-cost84} and add the attained-policy interval to obtain \eqref{eq:compatible-comparison84}. On the attained-cost event this entire argument is deterministic and holds for every valid collection of inequalities. Taking any number of allowances or any coefficient class therefore requires no additional probabilistic event. Finally,
$F(\theta)\le d+\epsilon\le\inf_\Theta F+\epsilon\le U+\epsilon$.
\end{proof}

A positive lower endpoint in \eqref{eq:compatible-comparison84} supports replacing $A$ by $B$ for every additional normalized charge below that endpoint. A strict interval contained in $(-m,m)$ gives equivalence at the declared margin $m$. These statements concern actual policy costs throughout the stated optimization classes. They do not identify a probability that an unrestricted fitting algorithm will enter those classes.

\subsection{Qualification with one common coefficient vector}
The fixed-label calibrated objective is
\[
 F(\theta)=\max_i(y_i^+-\phi_i^{\mathsf T}\theta)
           -\min_i(y_i^--\phi_i^{\mathsf T}\theta)
           +2\rho\bigl(G+\max_r|r_r^{\mathsf T}\theta|\bigr).
\]
A winner or joint-winner event adds linear inequalities
$d_j^{\mathsf T}\theta+c_j\le0$. The Numerical Appendix derives an exact lower certificate from two probability vectors, a derivative-budget multiplier and nonnegative event multipliers. All inequalities in a joint event use the same $\theta$ and the same qualification objective. A positive coefficient-box support separation certifies an empty event; a solver's infeasibility flag alone does not.

The complete-menu calculation and joint-event calculations can be stopped at any finite budget. Unchecked events add no exclusion, and inconclusive calculations retain the existing bound. Conversely, a valid new exclusion can immediately improve the economic assignment problem without any additional policy-cost observations. The arithmetic and integration costs required to obtain that exclusion remain separately chargeable.

\section{Complete-policy qualification and endogenous occupation}\label{sec:recursive86}
The first-date result holds the continuation policy fixed. A complete policy changes both later actions and the states at which they are taken. We now retain both changes. This extension concerns the original transition kernels on the continuous state space, rather than a replacement finite-state economy.

Fix a feasible reference policy $\tau$ and its bounded original-law values $V_t^\tau$. In the expectation model of Section~\ref{sec:core49}, define
\begin{equation}\label{eq:advantage86}
 A_t^\tau(x,a)=c_t(x,a)+\beta_tP_t^aV_{t+1}^\tau(x)-V_t^\tau(x).
\end{equation}
Here $a$ can denote a whole-cell feasible action rule, including acquisition and repair. At every date let $\mathcal C_t$ be a finite measurable acquisition partition. A menu label $a\in\mathcal A_{tC}$ denotes a measurable rule $x\mapsto a(x)\in A_t(x)$ for all $x\in C$. A parameter vector $\theta\in\Theta$ specifies the menu choice in every cell at every date; ties follow the declared measurable convention. The vector may share coordinates across dates and regimes. Let $\Theta_\epsilon=\{\theta\in\Theta:F(\theta)\le U+\epsilon\}$, where $U=F(\widehat\theta)$ for a feasible attained vector. The qualification functional may use the candidate's own future. Convexity is not required for the transfer; a claimed exclusion must nevertheless have a valid certificate for this actual functional.

Write $v=(t,C,a)$ for a menu event. Retain any family $\mathcal H_\epsilon$ of finite event sets for which a verified calculation proves
\begin{equation}\label{eq:conflict86}
 \Theta_\epsilon\cap\bigcap_{v\in K}\{\theta:\pi_\theta\text{ selects }v\}=\varnothing,
 \qquad K\in\mathcal H_\epsilon.
\end{equation}
A numerical infeasibility status alone is not such a certificate. The full-menu support duals in Section~\ref{sec:compatible84} give one implementation when labels and features yield the stated linear constraints. With nonlinear own-future dependence, those particular duals cannot simply be reused; verified outer relaxations or direct interval exclusion are required. Only strict separation excludes an event, so threshold equality remains included.

For every menu event and successor cell $D$, obtain uniform original-law enclosures
\begin{align}\label{eq:occupation-input86}
 \ell_{tCa}&\le A_t^\tau(x,a(x))\le u_{tCa}, &&x\in C,\notag\\
 \underline p_{tCaD}&\le P_t(D\mid x,a(x))\le\overline p_{tCaD}, &&x\in C.
\end{align}
The bounds must cover acquisition boundaries, feasible repairs and numerical error. Initial cell probabilities $\mu_C$ are known here; certified initial-probability constraints can instead be included in the feasible set. Averages under the reference policy are not substitutes for the uniform bounds in \eqref{eq:occupation-input86}.

Define $\mathcal P_\epsilon$ by nonnegative occupation variables $m_{tC}$, $q_{tCa}$, transition flows $z_{tCaD}$, and $0\le y_{tCa}\le1$, subject to
\begin{align}\label{eq:occupation-polytope86}
 m_{0C}&=\mu_C,& \sum_a q_{tCa}&=m_{tC},& \sum_a y_{tCa}&=1,\notag\\
 q_{tCa}&\le y_{tCa},& \sum_Dz_{tCaD}&=q_{tCa},&
 m_{t+1,D}&=\sum_{C,a}z_{tCaD},\notag\\
 \underline p_{tCaD}q_{tCa}&\le z_{tCaD}\le\overline p_{tCaD}q_{tCa},&&&\notag\\[-2pt]
 \sum_{v\in K}y_v&\le |K|-1 &&(K\in\mathcal H_\epsilon).&&
\end{align}
Every $\mu_C$ is nonnegative and $\sum_C\mu_C=1$. Thus each $m_t$ is a probability vector. These are outer occupation constraints, not a claim that the acquired cell process is Markov. In particular, they allow the within-cell state law to depend on all earlier actions.

\begin{theorem}[Recursive common-parameter cost transfer]\label{thm:occupation86}
Under the preceding boundedness, feasibility and enclosure conditions, put
\begin{equation}\label{eq:recursive-band86}
 L_\epsilon=\min_{\mathcal P_\epsilon}\sum_{t,C,a}B_t\ell_{tCa}q_{tCa},\qquad
 U_\epsilon=\max_{\mathcal P_\epsilon}\sum_{t,C,a}B_tu_{tCa}q_{tCa}.
\end{equation}
Every $\theta\in\Theta_\epsilon$ satisfies
\[
 L_\epsilon\le J_0^{\pi_\theta}-J_0^\tau\le U_\epsilon.
\]
The conclusion allows changes in every date's action and continuation, together with their induced original-law occupations. Valid additional compatibility or transition constraints can only tighten the enclosure. If $\epsilon$-dependent exclusions are retained monotonically, the enclosures are nested as qualification is strengthened. Certified dual lower and upper bounds suffice; an optimizer's reported optimum does not.

For two complete-policy classes with reference-cost interval $[l_{jk},u_{jk}]$ on a simultaneous event, their cost difference is enclosed on the same event by
\[
 [\,l_{jk}+L_j-U_k,\ u_{jk}+U_j-L_k\,].
\]
No probability union over the parameter vectors is added after that event. A qualified bound on a reference policy's loss against $V_0^*$ can also be enlarged by $U_\epsilon$, retaining the original optimum as the benchmark.
\end{theorem}

The proof in Supplement~\ref{supp:occupation86} uses a telescoping identity under the candidate policy's occupation law, then embeds each complete candidate in \eqref{eq:occupation-polytope86}. Independent choices at different cells and dates are not treated as jointly attainable. As in the first-date assignment bound, the relaxation itself can have unattainable points. The economic conclusion is an enclosure, not a worst-policy realization claim.

\subsection{Nonlinear recursion, temporal selves and games}\label{sec:nonlinear86}
For a monetary evaluation $\rho_t$ satisfying the monotonicity and cash-invariance conditions in Section~\ref{sec:core49}, let $\mathcal T_t^\pi$ be the fixed-policy operator. Suppose, uniformly over the states and policies in the declared class,
\[
 d_t^-\le\mathcal T_t^\pi V_{t+1}^\tau-V_t^\tau\le d_t^+.
\]
There is no linear-occupation identity for an arbitrary nonlinear $\rho_t$. Its correct counterpart is the following recursive order statement.

\begin{proposition}[Recursive order and deviation bounds]\label{prop:nonlinear86}
If $e_T^-=e_T^+=0$ and $e_t^\pm=d_t^\pm+\beta_te_{t+1}^\pm$, then
$e_t^-\le V_t^\pi-V_t^\tau\le e_t^+$ for every class member. More generally, on invariant recursive-value domains, if a monotone fixed-policy aggregator is $\lambda_t$-Lipschitz in its continuation argument and its reference residual is at most $b_t$ in absolute value, then
\[
 \|V_0^\pi-V_0^\tau\|_\infty
 \le\sum_{t<T}\left(\prod_{j<t}\lambda_j\right)b_t.
\]
For a dynamic game, apply the one-sided sandwich of Proposition~\ref{prop:sandwich49} to each player's best-response problem with the other players' complete strategies fixed. If that account is at most $r_{it}$ in every subgame, the attained profile has unilateral cost-improvement incentive at most $r_{it}$ there. This is an approximate equilibrium certificate, not an equilibrium-existence assertion. Date-specific temporal selves are treated as distinct decision makers, with their stated objectives and future strategies, rather than as one artificially time-consistent planner.
\end{proposition}

Endogenous preference states belong in the augmented state and transition before these hypotheses are checked. A diffusion approximation contributes its certified one-step operator error at every date, propagated with the same continuation moduli; an unproved weak-discretization rate is not assumed. The broader economic applications and their original hypotheses remain intact. The expectation identity, nonlinear order bound and strategic deviation certificate are distinct results, not interchangeable interpretations of one simulation table.

\section{Shared policy patterns and sharp occupation bounds}\label{sec:pattern89}
A common continuation restricts an entire contingent plan, including its prescriptions at states that are not reached. Pairwise exclusions need not retain that restriction after the occupation variables are relaxed. The following construction preserves the complete pattern before relaxing the state distribution. It strengthens Theorem~\ref{thm:occupation86}, rather than replacing the original Bellman problem with an unrelated finite-policy benchmark.

Use the original continuous-law model, reference policy, acquisition cells and whole-cell feasible action rules of Section~\ref{sec:recursive86}. A \emph{pattern} $\sigma$ assigns one menu label $\sigma(t,C)$ to every dated cell. Let $\Sigma_\epsilon$ be the patterns actually induced by members of $\Theta_\epsilon$, with the stipulated tie convention. This is a finite set even when the parameter set is infinite. For validity, it is sufficient to retain a finite superset $\overline\Sigma_\epsilon$ containing every actual pattern and satisfying all previously certified conflicts. A numerical search that finds some patterns does not establish such coverage. Conversely, asserting that a retained pattern is attained requires a qualified parameter witness; a fractional solution is not that witness.

For each retained pattern introduce a weight $\lambda_\sigma\ge0$, cell masses $m^\sigma_{tC}$ and transition masses $z^\sigma_{tCD}\ge0$. Define the disaggregated flow polytope by
\begin{align}\label{eq:pattern-flow89}
 \sum_\sigma\lambda_\sigma&=1,&
 m^\sigma_{0C}&=\lambda_\sigma\mu_C,\notag\\
 \sum_D z^\sigma_{tCD}&=m^\sigma_{tC},&
 m^\sigma_{t+1,D}&=\sum_C z^\sigma_{tCD},\notag\\
 \underline p_{tC\sigma(t,C)D}m^\sigma_{tC}
 &\le z^\sigma_{tCD}
 \le\overline p_{tC\sigma(t,C)D}m^\sigma_{tC}.&&
\end{align}
All initial probabilities, transition enclosures and advantage enclosures are the original-law inputs in \eqref{eq:occupation-input86}. In particular, the transition restrictions cover the whole cell and its feasible repair, not just an acquired point. Put
\begin{align}\label{eq:pattern-objective89}
 L_\epsilon^{\rm pat}
 &=\min\sum_{\sigma,t,C}B_t\ell_{tC\sigma(t,C)}m^\sigma_{tC},\notag\\
 U_\epsilon^{\rm pat}
 &=\max\sum_{\sigma,t,C}B_tu_{tC\sigma(t,C)}m^\sigma_{tC},
\end{align}
where both extrema are over \eqref{eq:pattern-flow89}.

\begin{theorem}[Pattern-consistent transfer and its approximation error]\label{thm:pattern89}
Suppose the qualification class is nonempty and the preceding coverage and enclosure conditions hold. Then every $\theta\in\Theta_\epsilon$ satisfies
\[
 L_\epsilon\le L_\epsilon^{\rm pat}
 \le J_0^{\pi_\theta}-J_0^\tau
 \le U_\epsilon^{\rm pat}\le U_\epsilon.
\]
The outer pair $(L_\epsilon,U_\epsilon)$ is that of Theorem~\ref{thm:occupation86} with the same conflicts and transition inputs. Valid exclusion of additional whole patterns can only improve the new pair. Certified lower and upper dual bounds remain sufficient for implementation.

Suppose further that the retained patterns are exactly the attainable patterns, and each has a qualification witness. Define
\begin{align*}
 \omega_t&=\max_{C,a}(u_{tCa}-\ell_{tCa}),&
 M_t&=\max_{C,a}\max\{|\ell_{tCa}|,|u_{tCa}|\},\\
 \delta_t&=\max_{C,a}\sum_D
       (\overline p_{tCaD}-\underline p_{tCaD}),&
 e_t&=\min\left\{2,\sum_{s<t}\delta_s\right\},
\end{align*}
with $e_0=0$, and put $\mathcal E=\sum_{t<T}B_t(\omega_t+M_te_t)$. If $j_-$ and $j_+$ are the minimum and maximum true class-cost differences, respectively, then
\begin{equation}\label{eq:pattern-error89}
 j_- -\mathcal E\le L_\epsilon^{\rm pat}\le j_-\le j_+
 \le U_\epsilon^{\rm pat}\le j_++\mathcal E.
\end{equation}
In particular, exact finite-state transitions and advantages give the exact class-cost extrema, attained by qualified patterns. More generally, any refinement sequence for which the displayed $\mathcal E$ tends to zero yields convergence of both bounds to the corresponding true extrema. The stated transition-width condition, not refinement alone, licenses this conclusion.
\end{theorem}

The full proof appears in Supplement~\ref{supp:pattern89}. It does not assume that the cell process of the continuous economy is Markov. Conditional cell-to-cell probabilities exist under each fixed policy at each date; the uniform enclosures contain those probabilities regardless of the within-cell distribution. Disaggregation prevents an optimization relaxation from buying one parameter pattern initially and a different pattern after observing its induced state.

\paragraph{Why complete patterns add information.}
The two-date example in Supplement~\ref{supp:occupation86} has two qualified signs of one shared coefficient. One sign visits the right-hand state and then selects its free action. The other sign selects the costly action there but never visits that state. Every actual policy therefore has cost zero. The pair-conflict flow relaxation permits mass $1/2$ on the right-hand state and its costly action; its upper bound is $1/2$. Disaggregation by the two attainable patterns gives upper bound zero. This strict improvement requires no new policy-cost sample, changed economic primitive or successful nonconvex fit.

\paragraph{Computation and scope.}
The number of complete patterns can be exponential in the number of dated cells. Enumeration is therefore a finite exact implementation, not a general polynomial-time claim. One may begin with an outer covering of parameter space, retain every unresolved pattern and refine only economically material branches. Every discarded branch needs a valid exclusion for the actual own-future qualification functional. Missing a branch invalidates the bound; retaining an unresolved branch only weakens it. The state-flow width in \eqref{eq:pattern-error89} and the global-pattern relaxation are distinct approximation costs. This distinction preserves both the original continuous economy and the original-optimum policy sandwich while making the added qualification work explicit.

\section{A common operator beyond the observed catalogue}\label{sec:population86}
A coefficient vector checked at finitely many tasks need not work at a new task. The constructive NBO has a stronger route when the economic primitives provide joint regularity. We state the necessary regularity on states and contexts rather than assuming successful extrapolation by a fitted nine-coefficient readout.

Let $r=(s,x)$ range over an invariant compact subset $K_t$ of $\R^{p_t}$, where $s$ contains declared economic contexts and any endogenous preference variables. Actions and transitions are defined on this augmented domain. Suppose the primitive moduli, Hausdorff-continuous feasible menus and measurable repairs in Theorem~\ref{thm:construct49} hold jointly in $r$, with constants $L_t,D_t,K_t^A$. Let an effective cover have $S_t$ nodes and radius $h_t$. At each node obtain its own-future minimized label with lower error $-e_t$ and upper error $D_tk_t+e_t$. Repair displacement is bounded by $K_t^A\|r-r_i\|_1+\zeta_t$.

\begin{theorem}[Population construction with one shared neural circuit]\label{thm:population86}
The single function
\[
 f_t(r)=\min_{i\le S_t}\{y_{ti}+L_t\|r-r_{ti}\|_1\}
\]
has an exact affine--ReLU circuit, a measurable original-owner actor and the bounds \eqref{eq:enclosures49} uniformly on $K_t$. Consequently \eqref{eq:constructionloss49}, with the joint-domain constants, holds for every initial context and every admissible evolution remaining in those domains. Contexts may be generated endogenously and adaptively. No independence assumption on their deployment sequence is needed for this uniform deterministic assertion.

The circuit uses $O(p_tS_t)$ nonzero scalar operations and weights, with absolute-value gates and a balanced minimum tree of depth $O(\log S_t)$. If the joint domains have effective $\ell^1$ covering numbers at most $C_th^{-k_t}$, this yields $S_t\le C_th_t^{-k_t}$. It is the explicitly shared joint-domain circuit, not a collection of task-specific networks. The original-owner and repair obligations remain necessary.
\end{theorem}

This is a constructive population theorem, not a claim that the small feature maps in the historical experiments satisfy its approximation bound. The covering dimension $k_t$, the ambient readout cost $p_t$, acquisition, integration and label work are charged separately. A low intrinsic covering dimension must be established for the declared economic family. The same cone also has a native min-plus implementation: exact equivalence of two encodings is not neural superiority.

\begin{proposition}[Acquisition coverage under persistent exploration]\label{prop:coverage86}
Choose $M$ deterministic radius-$h$ balls whose centers cover $K$. If, before each acquisition $i$, every such ball has conditional probability at least $a_h>0$ of containing the next node, then after $n$ acquisitions the probability that the acquired nodes fail to form a radius-$2h$ cover is at most $M(1-a_h)^n\le M\exp(-na_h)$. Thus $n\ge[\log M+\log(1/\delta)]/a_h$ is sufficient. The assertion permits an adaptive acquisition law satisfying the conditional lower bound; it does not silently replace such a law by IID sampling. On the coverage event, Theorem~\ref{thm:population86} applies with $h_t=2h$, provided the own-future labels and primitive moduli have their stated simultaneous certification.
\end{proposition}

\subsection{From approximation to predictor work}
The certificate's validity and a proposal's value are different. Consider a complete finite action menu at a query, with true objective $Q$ and estimates having simultaneous radius at most $e$. Select a minimum estimated objective and close only when its upper endpoint does not exceed every competitor's lower endpoint. Let $\Delta$ be the true best-to-second-best gap, with $\Delta=0$ for ties.

\begin{proposition}[A distributional work consequence of uniform accuracy]\label{prop:work86}
The preceding test closes whenever $\Delta>4e$. Suppose, uniformly over the relevant policy-induced query laws, the tie probability is at most $p_0$ and $\Pr(0<\Delta\le z)\le Cz^\alpha$ for $z>0$. If proposal and certification work are at most $h(e)$ and a protected fallback uses at most $W_+$, then
\[
 \E W\le h(e)+W_+\min\{1,p_0+C(4e)^\alpha\}.
\]
If the always-conventional comparator has expected work at least $W_-$, a positive saving follows when the displayed upper bound is below $W_-$. Comparing two upper bounds is not sufficient. The assertion is componentwise for nonnegative operation ledgers, or for a separately declared nonnegative resource-price vector. It is not a wall-clock assertion without a valid mapping from the ledger to time.
\end{proposition}

Uniformity over induced query laws is essential: a margin observed under a baseline state distribution does not establish the bound after the policy changes. The tie term cannot be discarded. A representation-specific empirical comparison must measure $h(e)$, achieved accuracy and the relevant query distribution for each representation. The constructed neural circuit provides an explicit member of this comparison; protected fallback alone supplies no evidence about neural training reliability. Proofs are in Supplement~\ref{supp:population86}.

\section{Original-law qualification under transport approximation}\label{sec:transport90}
The occupation bounds retain a complete policy and its induced state distribution. Their cellwise convergence result, however, requires shrinking transition-probability widths. Moving a deterministic successor by an arbitrarily small distance need not make those widths small. We develop a complementary certificate that tests the transition approximation against the actual stored continuation. It applies to singular as well as diffuse transition laws and does not require a continuous greedy actor.

Work in the original finite-horizon expected-cost model, with compact metric state spaces, nonempty feasible action correspondences, bounded costs and terminal cost $g$. All kernels below are probability kernels with finite first moments. Minima may be replaced by infima provided a measurable feasible approximate minimizer is supplied. Write $B_0=1$ and $B_t=\prod_{s<t}\beta_s$, where $\beta_s\ge0$. A parameter $\theta$ determines all stored continuations $f_{t,\theta}$ and a complete feasible policy $\pi_\theta$. Parameters, actions and continuations can be shared across dates and economic contexts; the state includes every endogenous variable needed by the original model.

Let $P_t(\cdot\mid x,a)$ denote the original transition law and $\widehat P_t(\cdot\mid x,a)$ a numerical approximation. The latter is an integration device, not a substituted economic primitive. Suppose that, uniformly on the declared domains and qualification class,
\begin{align}\label{eq:transport-input90}
 \operatorname{Lip}(f_{t+1,\theta})&\le K_{t+1},&
 W_1(P_t(\cdot\mid x,a),\widehat P_t(\cdot\mid x,a))&\le e_t,\notag\\
 \left|\widehat Q_{t,\theta}(x,a)-c_t(x,a)
 -\beta_t\widehat P_t^a f_{t+1,\theta}(x)\right|&\le\xi_t.&&
\end{align}
Here $W_1$ uses the stated metric on successor states, and $\widehat Q$ includes any additional certified quadrature and arithmetic approximation. In particular, $\xi_t$ is not set to zero merely because a computation returns a finite number. The bounds must hold for every original feasible action needed by the Bellman comparison, including any acquisition-boundary repair.

Define $\widehat T_{t,\theta}(x)=\inf_{a\in A_t(x)}\widehat Q_{t,\theta}(x,a)$. For each $\theta$ obtain constants $\ell_{t,\theta},u_{t,\theta}$ and $\gamma_{t,\theta}\ge0$ such that
\begin{align}\label{eq:transport-residual90}
 \ell_{t,\theta}&\le\widehat T_{t,\theta}(x)-f_{t,\theta}(x)
                    \le u_{t,\theta},\notag\\
 \widehat Q_{t,\theta}(x,\pi_{t,\theta}(x))
       -\widehat T_{t,\theta}(x)&\le\gamma_{t,\theta}
\end{align}
throughout the original state domain. A finite action search must include its certified gap to the full feasible menu in $\gamma$ or in the Bellman enclosures; a minimum over sampled actions alone is not the displayed infimum. Let $d^-_\theta\le g-f_{T,\theta}\le d^+_\theta$ uniformly, and put
\[
 b_t=\beta_tK_{t+1}e_t+\xi_t,\qquad
 \mathcal R_\theta=
 \sum_{t<T}B_t\{u_{t,\theta}-\ell_{t,\theta}
                  +\gamma_{t,\theta}+2b_t\}
       +B_T(d^+_\theta-d^-_\theta).
\]

\begin{theorem}[Transport-qualified original-optimum policy bound]\label{thm:transport90}
Under \eqref{eq:transport-input90}--\eqref{eq:transport-residual90}, every complete feasible candidate satisfies, at every initial state,
\begin{equation}\label{eq:transport-loss90}
 0\le V_0^{\pi_\theta}-V_0^*\le\mathcal R_\theta.
\end{equation}
Thus any common-parameter qualification class for which $\mathcal R_\theta\le\varepsilon$ is a class of complete policies with original-law loss at most $\varepsilon$. The result permits changes at every date and all corresponding changes in occupation. It requires neither a fixed continuation after the first date, independence of deployment contexts, nor Lipschitz continuity of the implemented policy.

The certificate is invariant to replacing each stored continuation by $f_{t,\theta}+k_{t,\theta}$, where the shifts are state independent, and shifting its approximate action values consistently. If the numerical transition is obtained by a measurable successor quantizer $q_t$ with $d(y,q_t(y))\le h_t$ almost surely, then $e_t=h_t$ is valid even when the two kernels are mutually singular. Its additional contribution is at most
\begin{equation}\label{eq:transport-quantization90}
 2\sum_{t<T}B_{t+1}K_{t+1}h_t.
\end{equation}
In particular, the transport part vanishes along any refinement for which this explicit sum tends to zero. No total-variation convergence is required. Convergence of the full certificate additionally requires the displayed residual, action-search, terminal and arithmetic contributions to vanish.
\end{theorem}

The full proof is in Supplement~\ref{supp:transport90}. It compares the candidate and the original Bellman optimum separately to the same stored continuation, before subtracting their bounds. Applying a baseline occupation law to a different complete policy is neither necessary nor valid. The use of a uniform continuation modulus is essential: a small transport distance alone does not control integration of an arbitrary discontinuous value function.

\paragraph{What the representation contributes.}
For the constructed affine--ReLU continuations, a valid Lipschitz modulus follows either from the cone construction or from verified layer norms. The resulting $K_{t+1}$ prices transition approximation in the policy certificate. The exact same function represented by a native minimum circuit has the same mathematical guarantee. The result identifies the role of the neural continuation without making a method-neutral verifier evidence of universal neural speed. Wasserstein robustness and finite-model approximation have established analyses \citep{zhou2024wasserstein,saldi2015finite}; our use is the uniform, translation-invariant qualification account for shared NBO continuations and their complete original-law policies.

\begin{corollary}[Discounted stationary qualification]\label{cor:stationary90}
In a stationary bounded-cost problem with $0\le\beta<1$, let $f$ be bounded and $K$-Lipschitz, and let a feasible stationary policy satisfy the preceding inequalities with constants $\ell,u,\gamma,e,\xi$. Then
\[
 0\le V^\pi-V^*
 \le\frac{u-\ell+\gamma+2(\beta K e+\xi)}{1-\beta}.
\]
This is an original-optimum certificate, not a claim of horizon-independent construction cost. The cost of acquiring a sufficiently accurate $f$, certifying $K$, integrating the original kernel and representing all relevant contexts remains part of the resource account.
\end{corollary}

\subsection{Joint allocation of residual and transition accuracy}
The coupled successor-cost program can now allocate two distinct sources of precision. If an implementation certifies residual-plus-greediness allowances $\eta_t$ and transport allowances $e_t$, its sufficient budget is
\begin{equation}\label{eq:joint-budget90}
 \sum_{t<T}B_t\eta_t+2\sum_{t<T}B_{t+1}K_{t+1}e_t
 +2\sum_{t<T}B_t\xi_t+B_T(d^+-d^-)\le\varepsilon.
\end{equation}
When fixed, verified continuation moduli and primitive resource estimates yield a positive-monomial upper model in $(\eta,e)$, Theorem~\ref{thm:coupled89} applies to its positive-weight coordinates. Coordinates with zero budget weight require separate caps and treatment; they do not satisfy that theorem's positive-weight premise. A modulus depending on an optimized parameter is not held fixed without a class-uniform upper certificate. This preserves the cost of transition precision rather than hiding it in the number of acquired states.

\subsection{Nonlinear and strategic evaluations}
For a cash-invariant monotone one-step risk evaluation, the same signed recursion applies when a separately established bound on the risk-evaluation discrepancy replaces $b_t$. Wasserstein distance and cash invariance alone do not imply such a bound. General monotone Lipschitz aggregators use the absolute-error recursion of Proposition~\ref{prop:nonlinear86}, with their own continuation moduli, rather than an unsupported linear expectation identity. For games, the transport account enters each fixed-opponent best-response certificate separately. Preference states and temporal selves retain the augmented-state and unilateral-deviation interpretation already stated in Section~\ref{sec:nonlinear86}. The new result supplements, rather than removes, the original diffusion, recursive-preference and game constructions.

\paragraph{A singular-transition diagnostic.}
Take one decision date, $\beta=1$, terminal cost $g(x)=x$ on $[0,1]$, and actions $0$ and $1$. Action $0$ has zero immediate cost and true successor $h>0$; action $1$ costs $h/2$ and has successor zero. The numerical kernel places both successors at zero. With $f_0=0$ and $f_1=g$, its greedy choice is action $0$, its residual span is zero and its true regret is $h/2$. Omitting the transition charge would falsely certify zero regret. The transport certificate is $2h$ and converges to zero, while the total-variation distance for action $0$ is one for every $h>0$. Appendix~\ref{app:transport90} gives exact rational checks. This diagnostic is not a new observation in the investment experiment.

\section{Horizon complexity and the price of accuracy}\label{sec:complexity86}
The inherited recursive query experiment exposes descendant growth. The own-future constructive backend has a different dependency graph: once $f_{t+1}$ is stored, a query evaluates that representation and does not solve the remaining Bellman problem again. This distinction changes the horizon account without removing any of the recursive benchmarks.

For each date let $S_t$ be its state-context cover size, $J_{ti}$ its feasible action-net size at node $i$, $M_t$ its certified integration work in primitive representatives, and $E_{t+1}$ the cost of evaluating its stored successor. Construction uses
\begin{equation}\label{eq:dagwork86}
 Q_t=M_t\sum_{i\le S_t}J_{ti},\qquad
 W\le\sum_{t<T}\{Q_t(C_t^{\rm prim}+E_{t+1})+C_t^{\rm compile}\}
       +C_T^{\rm terminal}+C^{\rm output}.
\end{equation}
The formula counts a directed acyclic sequence of stored continuations. A literal descendant solver instead multiplies successor solves by its action and shock branching factors. Neither expression is a lower bound for every possible algorithm. The parameters in \eqref{eq:dagwork86} may themselves grow with horizon if the primitive regularity deteriorates.

\begin{proposition}[Optimal dated error allocation]\label{prop:allocation86}
Suppose a certified constructive implementation has a resource upper account $a_t\eta_t^{-p}$ for attaining total date-$t$ residual allowance $\eta_t>0$, with common $p>0$. Each $a_t$ includes action search, integration, feasibility, compilation and the cost of evaluating the stored successor; it cannot hide another recursive Bellman solve. Put $w_t=B_t>0$ and
\[
 S=\sum_t a_t^{1/(p+1)}w_t^{p/(p+1)}.
\]
Among allocations satisfying $\sum_tw_t\eta_t\le\varepsilon$, the minimum of this resource upper model and its minimizer are
\begin{equation}\label{eq:allocation86}
 W_{\rm model}^*=S^{p+1}\varepsilon^{-p},\qquad
 \eta_t^*=\frac{\varepsilon}{S}(a_t/w_t)^{1/(p+1)}.
\end{equation}
This optimizes the declared upper model, not the unknown fastest algorithm. Rounding each allowance down to a dyadic number within a factor two preserves the error constraint and increases that model by at most $2^p$.

If $a_t\le A$ and $w_t=\beta^t$ for $0<\beta<1$, then
\[
 W_{\rm model}^*\le
 A\varepsilon^{-p}(1-\beta^{p/(p+1)})^{-(p+1)}.
\]
With a maximum permissible allowance $\eta_{\max}$ and positive per-date setup costs, clipping the displayed allocation adds at most
$\sum_t a_t\eta_{\max}^{-p}$ to this upper bound, plus the explicit setup account. For $w_t=1$, the un-clipped bound is $A(T+1)^{p+1}\varepsilon^{-p}$. Thus uniform regularity and reusable stored continuations give polynomial, rather than descendant-tree, horizon dependence at fixed dimensions. They do not give dimension-free computation.
\end{proposition}

The power model follows, for example, when effective state, action and innovation covers and their certified errors have polynomial work bounds, successor evaluation has the separately proved compiled bound, and all moduli are uniformly controlled. The exponent must include all these factors. Tensor compilation in Theorem~\ref{thm:compile49} has a $2^{p_t}$ query factor and cell-location work; a dense neural cone has an $S_{t+1}$ factor. Those are different implementation accounts and must not be interchanged.

\subsection{Operation vectors, arithmetic and reliability}\label{sec:ledger86}
We retain the resource vector
\[
 (N_{\rm primitive},N_{\rm descendant},N_{\rm interval},N_{\rm rational},
 N_{\rm network},B_{\rm read},B_{\rm write},M_{\rm peak}).
\]
Counts are reported before scalar prices are applied. Rational operations additionally record operand bit lengths; an arithmetic-operation count is not a bit-complexity bound. Multiplication and division costs at those lengths enter a bit-cost account. Loading, validation, serialization and shutdown belong to complete process time. Peak memory is a capacity requirement, not an additive flow cost unless its duration or rental convention is specified.

A fitting warning, invalid proposal, zero-progress probe, successful screen, protected fallback and terminal service status are different events. In particular, zero new fits cannot estimate fitting reliability, and a successful fallback cannot be counted as a successful prediction. The retained 336-service recursive catalogue and the nine-controller attained-cost experiment have different units of observation. New mathematical fixtures below establish the operation identities and allocation formula; they are not relabeled as a re-execution of the inherited eight-million-query vector tree or as a hardware calibration. The complete proof of Proposition~\ref{prop:allocation86} is in Supplement~\ref{supp:allocation86}.

\section{Accuracy allocation with endogenous successor cost}\label{sec:coupled89}
The separable allocation in Proposition~\ref{prop:allocation86} applies when its coefficients are fixed independently of the allowances being optimized. A dense successor circuit generally has a size that depends on its own allowance. Charging that size at the preceding date couples the two choices. We now optimize this dependency explicitly.

Let $\eta=(\eta_0,\ldots,\eta_{T-1})$ be positive dated residual allowances. Primitive analysis of a specified stored-continuation implementation supplies the resource upper model
\begin{equation}\label{eq:coupled-work89}
 F(\eta)=F_0+\sum_{j=1}^{J}c_j\prod_{t<T}\eta_t^{-a_{jt}},
 \qquad c_j>0,\quad a_{jt}\ge0.
\end{equation}
The coefficients, exponents and setup term $F_0\ge0$ are fixed before choosing $\eta$. A term $c_t\eta_t^{-p_t}\eta_{t+1}^{-s_{t+1}}$ charges date-$t$ label acquisition and evaluation of its stored successor. Compilation, terminal evaluation and storage-related work may supply other terms. An operation ledger is scalarized only at declared nonnegative resource prices; the model does not estimate monetary prices or energy use.

The feasible accuracy set is
\begin{equation}\label{eq:coupled-budget89}
 \mathcal A=\{\eta:0<\eta_t\le\bar\eta_t,\quad
                    \sum_{t<T}w_t\eta_t\le\varepsilon\},
 \qquad w_t>0,\quad\bar\eta_t>0,\quad\varepsilon>0.
\end{equation}
For the original expected-cost sandwich, $w_t=B_t$. The caps specify where the implementation's error and resource bounds hold. A coordinate with no resource-dependent term can be fixed separately; assume here that every coordinate has a positive exponent in at least one term.

\begin{theorem}[Coupled accuracy optimization with an exact certificate]\label{thm:coupled89}
The upper-model minimization $F_* =\min_{\mathcal A}F(\eta)$ has a unique positive minimizer. It is a convex program in $\eta$ and, after $\eta_t=\exp(x_t)$, also has a convex exponential representation. Standard geometric-programming methods can propose an allocation, but their numerical status is not its certificate.

For any feasible $\eta^0$, write $g=\nabla F(\eta^0)$. Let $\lambda\ge0$ and $\nu_t\ge0$ satisfy
\begin{equation}\label{eq:coupled-certificate89}
 g_t+\lambda w_t+\nu_t\ge0\qquad(t<T).
\end{equation}
Then the independently checkable bounds are
\begin{align}\label{eq:coupled-dual89}
 D(\eta^0,\lambda,\nu)
 &=F(\eta^0)-g\cdot\eta^0-\lambda\varepsilon
                  -\sum_t\nu_t\bar\eta_t,\notag\\
 D(\eta^0,\lambda,\nu)&\le F_*\le F(\eta^0).
\end{align}
At the optimizer, the Karush--Kuhn--Tucker multipliers make these bounds equal. For any candidate, the choice $\nu=0$ and
$\lambda=\max_t\{-g_t/w_t\}$ supplies a valid, possibly loose, lower bound. With integer exponents and rational inputs and candidates, every check is rational arithmetic. For other exponents, certified enclosures of the function and derivatives are required.

Put $P=\max_j\sum_ta_{jt}$. Rounding each allowance down by less than a factor two preserves feasibility and gives
\[
 F(\eta^{\rm dyad})-F_0
 \le2^P\{F(\eta^0)-F_0\}.
\]
The exponent is the total degree of a coupled term, not merely the exponent of the date at which that term is paid.
\end{theorem}

The proof in Supplement~\ref{supp:coupled89} includes attainment, strict convexity and the certificate. Convex duality and geometric programming are established techniques \citep{boyd2004,boyd2007gp}; their role here is to preserve the successor dependency in the NBO's original accuracy account. The certificate bounds the optimum of an explicitly stated resource upper model. It is not a lower bound on the work of every algorithm solving the economic problem.

\begin{corollary}[A constructive horizon bound with coupled evaluation]\label{cor:coupled-horizon89}
Suppose $w_t=1$, $\varepsilon/T\le\min\{1,\min_t\bar\eta_t\}$, $J\le KT$, $c_j\le C$, and the total degrees are bounded by a fixed $P$. Then a feasible stored-continuation implementation satisfies
\[
 F_*\le F_0+KC\,T^{P+1}\varepsilon^{-P}.
\]
For a dense successor account with label exponent $p$ and evaluation exponent $s$, the adjacent-date term has degree $p+s$. Thus a valid polynomial horizon account charges $p+s$, rather than silently holding the successor cost constant. Worsening regularity, growing cover dimension, bit lengths and setup remain explicit inputs; none is bounded by this corollary without its stated hypotheses.
\end{corollary}

\paragraph{A priced two-date allocation.}
Consider the exact upper model
\[
 F(\eta_0,\eta_1)=\eta_0^{-1}
             +(\eta_0\eta_1)^{-1}+2\eta_1^{-2},
 \qquad\eta_0+\eta_1\le\tfrac12,\quad\bar\eta_t=1.
\]
The equal allocation costs $52$. The current numerical appendix gives a rational feasible allocation, its independently recomputed lower certificate and its gap. The example isolates the accuracy dependency algebraically; its resource units are not fitted runtime measurements. It also provides a regression test that rejects a negative multiplier, an infeasible budget or a forged lower bound. In the original economy the same construction applies only after its coefficients have been derived from the actual label, action, integration and stored-circuit accounts.

\paragraph{Policy and resource certificates.}
Qualification determines which policy outputs meet the economic accuracy requirement. The coupled program determines how a specified implementation distributes the corresponding dated allowances. These are different optimizations. A tight resource certificate cannot replace an original-law policy certificate, and a tight policy certificate does not determine hardware time. Keeping both accounts explicit permits constructive improvement without a claim of dimension-free or universally faster computation.

\section{Sequential qualification of policy costs}\label{sec:sequential73}
A fixed-sample interval does not by itself justify repeated inspection until a controller appears acceptable. Nor do cost observations from changing controllers identify the last controller's cost. We extend the frozen-continuation identity to predictable episodes and distinguish these targets. This uses established time-uniform concentration principles \citep{howard2020time}; confidence sequences have also been used for gated deployment in contextual bandits \citep{karampatziakis2021}. Here the observations are verified, continuous-law NBO cost contrasts, not importance-weighted bandit rewards.

Let $\mathcal F_{i-1}$ contain all information before episode $i$. Conditional on it, choose two complete controllers, a common initial law, a continuation $h_i$ fixed through the episode, and valid coefficients of Theorem~\ref{thm:frozen72}. The resulting cost contrast $D_i$ has conditional mean $\Delta_i$ equal to the original-law expected cost difference and predictable support $[a_i,b_i]$. Let $[L_i,U_i]\subseteq[a_i,b_i]$ enclose $D_i$ for the entire continuous innovation cell. Write $V_n=\sum_{i\le n}(b_i-a_i)^2$. For fixed positive numbers $\lambda_j,w_j$, $1\le j\le J$, with $\sum_jw_j\le1$, define
\begin{equation}\label{eq:anytime-boundary73}
 B_n=\min_{1\le j\le J}
       \left\{\frac{\log(2/(\alpha w_j))}{\lambda_j}
                    +\frac{\lambda_j V_n}{8}\right\}.
\end{equation}

\begin{theorem}[Predictable original-law cost certificate]\label{thm:anytime73}
Under the preceding conditional-law, support and enclosure conditions, with probability at least $1-\alpha$, simultaneously for every $n\ge1$,
\begin{equation}\label{eq:anytime-interval73}
 \frac{\sum_{i\le n}L_i-B_n}{n}
 \ \le\ \frac1n\sum_{i\le n}\Delta_i
 \ \le\ \frac{\sum_{i\le n}U_i+B_n}{n}.
\end{equation}
Each endpoint may also be intersected with the corresponding average predictable support endpoint. The assertion holds at every almost surely finite stopping time. If the complete controller pair and initial-law/reset contract stay fixed, then $\Delta_i=\Delta$, even when the qualified continuation changes predictably between episodes. Otherwise the target is the average conditional cost difference, not the last controller's cost.
\end{theorem}

The proof applies conditional Hoeffding bounds to each endpoint and Ville's inequality to a finite collection of nonnegative supermartingales. It requires no independent or identically distributed deployment contexts. It does require the stated conditional transition laws. Within-episode cache adaptation is allowed only when the controllers are treated as history-dependent algorithms and the greedy-gap qualification covers every memory state that can occur. Repeating the same program identifier while carrying a different cache state does not establish a fixed conditional target.

The implementation uses a finite dyadic collection of $\lambda_j$ and exact rational upper bounds for the logarithms. These give conservative, computable versions of \eqref{eq:anytime-boundary73}. It reserves the support and target identity before accepting an episode. A failed numerical evaluation consumes the full previously qualified support interval and its actual work; it is not deleted after inspecting its outcome. A failure of the support certificate itself stops qualification.

\begin{corollary}[Costed replacement after optional stopping]\label{cor:adoption73}
For a fixed target pair, let $[l_N,u_N]$ be the interval at a finite stopping time $N$. Let a declared deployment count be $m\ge1$, a one-time replacement and qualification charge be $F\ge0$, and the incremental per-deployment charge be $c$, all in the same units as the policy cost. On the simultaneous event of Theorem~\ref{thm:anytime73}, the net incremental cost lies in
\[
 [l_N+F/m+c,\ u_N+F/m+c].
\]
Consequently $u_N+F/m+c<0$ certifies strict net improvement. Inclusion of $[l_N,u_N]$ in a declared equivalence band certifies equivalence at that band, not strict superiority. New adaptively selected controller pairs may be qualified in separate epochs with conditional error allocations whose sum is at most the overall error budget.
\end{corollary}

A wall-clock reading is not a monetary charge. Measured setup, validation, rejected evaluations and successful retrievals are retained separately; prices are explicit decision inputs. The number of qualification episodes is not automatically the number of deployments over which a setup cost is amortized. These distinctions let the original investment decision use a valid stopping rule without silently changing either its estimand or its units.

\noindent The following attained-cost experiment retains the legacy coefficientwise qualification rule. Section~\ref{sec:study81} subsequently tests calibrated qualification on a separately frozen factorial catalogue.
\section{Original-law qualification of constructed policies}\label{sec:cost77}
A continuation can have a verified fitting objective without its implemented policy having an identified economic advantage. The distinction is resolved by evaluating the same complete policy that the construction returns. We now connect the qualification-optimal readouts of Section~\ref{sec:qualification73} to independent original-law cost observations. Neither the fitted future nor the acquired actor is replaced during evaluation. The comparison concerns the constructed NBO policies, rather than a different recursive controller supplied with the same neural label.

\subsection{Conditional integration and numerical observations}
In the two-state, two-date investment economy, write $a_t=\pi_t(X_t)$ and $\beta=15/16$. A cash observation and a conditionally integrated observation are
\begin{align}
 C^\pi&=c(X_0,a_0)+\beta c(X_1,a_1)+\beta^2g(X_2),\label{eq:cash77}\\
 Y^\pi&=c(X_0,a_0)+\beta c(X_1,a_1)
                    +\beta^2\bar g(X_1,a_1),\label{eq:conditional77}
\end{align}
where $\bar g(x,a)=P^ag(x)$. If $m(x,a)$ is the deterministic part of the transition, the original common, opposite-signed innovation gives the exact identity
\begin{equation}\label{eq:last-moment77}
 \bar g(x,a)=g(m(x,a))+\frac{5}{3072}.
\end{equation}
The shortage term depends on the sum of the coordinates and is unchanged by that innovation. The quadratic state and difference terms supply the variance correction. Consequently,
\[
 \E Y^\pi=\E C^\pi=\E J_0^\pi(X_0),\qquad
 \operatorname{Var}(Y^\pi-Y^\sigma)
       \le\operatorname{Var}(C^\pi-C^\sigma)
\]
when both policies use the same initial state and innovations. The variance inequality concerns the exact observations; it does not assert that every finite-sample numerical interval must be narrower.

The state cube is invariant. On that cube, $0\le a\le1/4$, a stage cost is at most $103/64$, and the terminal cost is at most $37/16$. Both observations therefore lie in
\begin{equation}\label{eq:support77}
 [0,H],\qquad H=\frac{103}{64}(1+\beta)+\frac{37}{16}\beta^2
                  =\frac{21097}{4096}.
\end{equation}
A paired difference has the signed support $[-H,H]$. Negative common-path gains are not clipped. No original-optimum value query is required to compute either observation for a frozen acquired actor.

An initial or innovation bin is a continuous cell, not a point mass. The evaluator encloses every true state in the bin, every acquired cell reached by that state enclosure, and every resulting feasible action. When an enclosure crosses an acquisition boundary, all intersected actions enter the continuation enclosure. A second implementation evaluates the same primitives using integer-directed arithmetic and expanded polynomials. The reported endpoints are the hull of the two valid enclosures, intersected with the proved support. Arithmetic agreement is a check on the implementations; validity follows from the enclosure rules and the primitive identities, not from agreement alone.

\subsection{Precision and the decision criterion}
The following statement applies to any fixed finite collection of qualified policies, not only to the two feature bases used below. All construction information is fixed before the independent evaluation rows are drawn. Dependence among policies on a common row is allowed.

\begin{theorem}[Original-law cost qualification]\label{thm:precision77}
For each of $K$ declared estimands, let $Z_i$ be its original-law observation on independent row $i$, and suppose a deterministic full-bin calculation returns
\[
 a\le L_i\le Z_i\le U_i\le b,\qquad i=1,\ldots,N.
\]
The inclusion may instead hold for the original conditional expectation within the sampled bin. Put $R=b-a$, let $\bar L,\bar U$ denote sample means, and let $s_L^2,s_U^2$ be the unbiased endpoint sample variances. For $N\ge2$ and $\ell>0$, define
\[
 r_L=\sqrt{\frac{2s_L^2\ell}{N}}+\frac{7R\ell}{3(N-1)},\qquad
 r_U=\sqrt{\frac{2s_U^2\ell}{N}}+\frac{7R\ell}{3(N-1)}.
\]
With probability at least $1-4K e^{-\ell}$, simultaneously for the declared collection,
\begin{equation}\label{eq:interval77}
 \E Z_i\in[\bar L-r_L,\bar U+r_U].
\end{equation}
Before rounding the final endpoints, its width is exactly
\begin{equation}\label{eq:width77}
 (\bar U-\bar L)
 +\sqrt{\frac{2s_L^2\ell}{N}}+\sqrt{\frac{2s_U^2\ell}{N}}
 +\frac{14R\ell}{3(N-1)}.
\end{equation}
Upward square roots and outward final rounding preserve coverage. The first term includes all retained numerical and acquisition uncertainty; the middle terms reflect observed endpoint dispersion; the last term is the finite-support penalty. A declared full-width target is met only when their computed sum meets it.
\end{theorem}

The proof in the Technical Supplement applies the empirical-Bernstein inequality of \citet{maurer2009} to the lower and upper endpoint variables. The statistical inequality is established work. The contribution here is its application to the original-law costs of the actual constructed actors, with executable whole-bin observations, independent arithmetic and an explicit precision account. Refining only a simulation bin cannot remove the finite-support or sampling terms. Conversely, an inaccurate uniform regret certificate need not prevent an accurate comparison of two actual policy costs.

Precision, equivalence and improvement are different conclusions. For a cost-difference interval $[l,u]$ and a predeclared margin $\delta$, equivalence requires $-\delta<l\le u<\delta$. Noninferiority of the first policy requires $u<\delta$; strict superiority requires $u<0$. An interval can meet its width target and fail equivalence because it identifies a material difference. These decisions are made from exact endpoints, not from rounded table entries.

\begin{corollary}[Priced selection and adoption]\label{cor:purchase77}
On the simultaneous event of Theorem~\ref{thm:precision77}, let $[L_j,U_j]$ qualify the cost of policy $j$, and let $w_j$ be its declared nonnegative resource charge. A minimizer $\widehat j$ of $U_j+w_j$ satisfies, for every policy $j$,
\[
 J_{\widehat j}+w_{\widehat j}
       \le J_j+w_j+(U_j-L_j).
\]
The inequality holds for every assignment of resource prices on the same event. If the gain from replacing an installed controller is in $[l,u]$, then $m$ deployments, monetary scale $v\ge0$, qualification charge $q\ge0$ and replacement fee $f\ge0$ give expected net benefit in
\[
 [mv l-q-f,\;mv u-q-f].
\]
In particular, when $l>0$ and $v>0$, every $m>(q+f)/(vl)$ is certified beneficial.
\end{corollary}
\begin{proof}
The first assertion follows from $J_{\widehat j}\le U_{\widehat j}$, the minimizing property, and $U_j\le J_j+U_j-L_j$. These deterministic inequalities hold for all nonnegative prices once all cost intervals cover. The second assertion is multiplication by $mv$ and subtraction of $q+f$.
\end{proof}

The resource charges in this result must be specified, not inferred from normalized losses. An already paid construction and a prospective purchase have different charges. When only a shared qualification service has been timed, charging that entire service to a proposed replacement is conservative; inventing a per-policy marginal qualification time is not.

\subsection{The constructed-controller experiment}
The fixed catalogue consists of the four previously constructed NBO actors and the admissible zero-investment incumbent. ReLU and degree-four Bernstein continuations each have nine coefficients and are constructed on eight- or sixteen-cell coordinate covers. At a fixed cover, the two bases use the same acquisition contract and primitive information. The actor arrays, actual fitted futures and exact optimization certificates are unchanged. The zero-investment policy supplies an economically interpretable installed reference; it is not substituted for the strong conventional query competitors in the preceding study.

Before execution, the protocol fixes $N=524{,}288$, forty-bit continuous bins, the sample stream, the full-width target $1/256$, and the separate two-sided equivalence margin $1/256$. All five marginal costs and ten paired differences are reported for both observations in \eqref{eq:cash77}--\eqref{eq:conditional77}. Thus $K=30$ and $\ell=10$ give a joint failure bound below $1/100$. Reusing a row for another controller or estimator does not increase the independent sample count. The declared independence model concerns the simulated rows; a fixed pseudorandom seed establishes reproducibility, not mathematical independence.

\begin{table}[htbp]\centering
\caption{Original-law costs of the constructed controllers}\label{tab:cost77}
{\small\setlength{\tabcolsep}{4pt}\begin{tabular}{lcc}
\toprule
Controller & Expected cost & Gain over zero incumbent \\
\midrule
Zero & [0.773561, 0.780120] & -- \\
R8 & [0.607716, 0.613107] & [0.165305, 0.167552] \\
B8 & [0.684034, 0.690159] & [0.089044, 0.090443] \\
R16 & [0.607022, 0.612384] & [0.165995, 0.168280] \\
B16 & [0.682910, 0.689022] & [0.090168, 0.091580] \\
\bottomrule\end{tabular}}
\par\smallskip{\footnotesize R and B denote the original nine-coefficient ReLU and Bernstein constructions; the number is the coordinate cover size. All policies are frozen before the new evaluation rows. The displayed observations integrate the last innovation analytically without changing expected cost. Endpoints are rounded outward. Paired gain intervals, not differences of marginal cost intervals, determine gain signs.}
\end{table}

\begin{table}[htbp]\centering
\caption{Paired feature and acquisition comparisons}\label{tab:pairs77}
{\small\setlength{\tabcolsep}{4pt}\begin{tabular}{lcrc}
\toprule
First minus second & Expected-cost difference & Full width & Equivalent \\
\midrule
R8 $-$ B8 & [-0.077568, -0.075801] & 0.001767 & No \\
R16 $-$ B16 & [-0.077159, -0.075368] & 0.001791 & No \\
R8 $-$ R16 & [0.000227, 0.001191] & 0.000963 & Yes \\
B8 $-$ B16 & [0.000661, 0.001599] & 0.000937 & Yes \\
\bottomrule\end{tabular}}
\par\smallskip{\footnotesize A negative interval identifies lower cost for the first controller. Equivalence is determined from exact endpoints at the predeclared two-sided margin 1/256. The separate full-width target is also 1/256. Matched-basis comparisons hold the acquisition contract fixed; cross-resolution comparisons do not.}
\end{table}

\paragraph{Outcomes of the fixed precision design.}

Across both estimators, 20 of the twenty paired intervals meet the declared full-width target. Their widths range from 0.000936 to 0.002285. Exactly 4 meet the separate equivalence criterion. Neither a failed width target nor a failed equivalence comparison is omitted.

The conditionally integrated comparisons identify a positive gain over the zero-investment incumbent for 4 of the four constructed policies. These are gains of complete implemented controllers, not reductions in a qualification objective.

At the two matched acquisition resolutions, the ReLU-minus-Bernstein comparisons identify 2 lower-cost ReLU outcomes, 0 higher-cost outcomes, and 0 unresolved signs. This finite comparison is not a universal representation ranking.


The construction and the evaluation now refer to the same complete controllers. The resulting cost intervals, rather than a reduction in the residual-fitting objective, determine the adoption regions in Corollary~\ref{cor:purchase77}. Comparisons between the feature bases identify only what their paired intervals support. The earlier unresolved and adverse query-method comparisons remain distinct evidence and are preserved in the Numerical and Development Appendix.

Both interval implementations evaluate every new trajectory and every acquired-action rectangle. Additional rational midpoint checks test predetermined rows. The phase ledger separates controller loading, lookup preparation, binary64 evaluation, integer evaluation, exact statistical arithmetic, output and the complete parent process. The independent replay recomputes the entire new enclosure catalogue without replacing its original clocks. This evaluates the two-date constructed-policy problem; it is not an independent rerun of the older, much larger vector Bellman tree.

\section{Qualification, representation, and attained economic cost}\label{sec:study81}
The preceding attained-controller experiment used a coefficientwise qualification bound. Its ReLU--Bernstein ordering is an observation about those specified constructions. We now examine whether that ordering persists after the two bases receive the same qualification charge on their common affine functions. This directly tests a source of representation-dependent conservatism while preserving the original nonlinear investment economy and the Neural Bellman Operator construction.

\subsection{A prospective factorial comparison}
The design fixes two bases, two qualification rules and two acquisition resolutions. The eight constructed controllers and the installed zero controller are evaluated on the same original-law paths. Each basis has nine coefficients with absolute bounds eight, the same node covers, thirty-three-action search menus and downward action quantum $1/4096$. The calibrated rule replaces the coefficientwise modulus by \eqref{eq:relu-modulus81} or \eqref{eq:bernstein-modulus81}, both in residual qualification and in the regularity allowance passed to the previous fitted date. Both own-future continuations are constructed backwards. An exact primal--dual gap of at most $1/65536$ qualifies each new fit. The coefficient proposal is quantized at forty bits; no hidden-feature optimizer or outcome-selected basis is introduced.

A source manifest was committed before construction, and the complete actor identities were committed before generating 524,288 new four-coordinate continuous-bin rows. The bin precision is forty bits. Two estimators, cash cost and exact last-innovation conditional expectation, evaluate the same expected costs under the original uniform initial law and continuous innovations. For each estimator the frozen family contains nine marginal costs, thirty-six paired differences and two basis-by-qualification interactions. All ninety-four estimands share a single failure allowance bounded by $376\exp(-11)<1/100$. Shared paths and the two estimators do not multiply the number of independent rows.

The whole-bin enclosures are independently computed in outward binary64 arithmetic and directed-integer arithmetic. Their hull, intersected with the proved support, is used for inference. An exact rational implementation additionally checks 512 midpoint trajectories per controller. All nine controllers, bins, endpoints, exact sample moments and disagreements are part of the retained protocol. This revision re-evaluates every bin and every inference quantity; it neither regenerates those observations as a second sample nor substitutes replay time for the original service clock.

\begin{table}[htbp]
\centering
\caption{Basis comparison before and after qualification calibration}\label{tab:basis81}
{\small\setlength{\tabcolsep}{3.5pt}
\begin{tabular}{lrcc}
\toprule
Qualification & $n$ & ReLU minus Bernstein & Lower cost \\
\midrule
Legacy & 8 & [-0.077633, -0.075733] & R \\
Legacy & 16 & [-0.077232, -0.075307] & R \\
Calibrated & 8 & [0.003571, 0.004688] & B \\
Calibrated & 16 & [0.000955, 0.002013] & B \\
\bottomrule\end{tabular}}
\par\smallskip{\footnotesize All entries are simultaneous paired original-law cost intervals from the same prospective R80 family, with exact last-innovation conditional integration. Positive values favor Bernstein. Equal coefficient counts and affine-exact modulus charging do not make the two nonlinear function classes identical. The old fits are evaluated, not overwritten.}
\end{table}

\begin{table}[htbp]
\centering
\caption{Actual controller gains from calibrated qualification}\label{tab:calibration-gains81}
{\small\setlength{\tabcolsep}{3.5pt}
\begin{tabular}{lrc}
\toprule
Basis & $n$ & Legacy minus calibrated cost \\
\midrule
Relu & 8 & [0.000457, 0.001653] \\
Relu & 16 & [0.005159, 0.006341] \\
Bernstein & 8 & [0.080964, 0.082772] \\
Bernstein & 16 & [0.082601, 0.084405] \\
\bottomrule\end{tabular}}
\par\smallskip{\footnotesize Each comparison retains its basis, coefficient count and bounds, robust action menus, observation resolution and original economic primitives. Both attained future fits are reconstructed under the stated qualification rule. The displayed gains are not timing estimates or claims about unrestricted neural fitting.}
\end{table}


The calibration improves the expected cost of each basis at each resolution. In particular, the fine ReLU controller gains between $0.005159$ and $0.006341$ relative to its legacy construction. The matched basis ordering nevertheless reverses: calibrated Bernstein has lower cost than calibrated ReLU at both resolutions. At $n=16$ the ReLU-minus-Bernstein interval is $[0.000955,0.002013]$, rounded outward. That strict ordering coexists with equivalence at the declared margin $1/256$; a small positive difference can satisfy both statements. The corresponding coarse-resolution interval does not establish equivalence at that margin.

\begin{table}[htbp]
\centering
\caption{Predeclared basis-by-qualification interactions}\label{tab:interactions81}
{\small\setlength{\tabcolsep}{3.5pt}
\begin{tabular}{rc}
\toprule
$n$ & Change in Bernstein-minus-ReLU cost \\
\midrule
8 & [-0.082254, -0.079371] \\
16 & [-0.079209, -0.076297] \\
\bottomrule\end{tabular}}
\par\smallskip{\footnotesize The interaction is (calibrated Bernstein minus calibrated ReLU) minus (legacy Bernstein minus legacy ReLU). These four-term contrasts were included in the source-frozen family before the new cost rows were generated. They are not obtained by treating the four marginal estimates as independent.}
\end{table}


The predeclared interactions isolate the change in the basis comparison. They do not establish that the bases span identical functions, that the Bernstein derivative-control bound is globally exact, or that every other numerical difference has been removed. The positive operational conclusion is more specific: charging the function's derivative structure changes the attained original-law policy cost, and a policy ranking cannot be attributed to a representation name without examining its qualification rule. The historical ReLU gain and the new reversed matched ranking remain separate valid observations.

\subsection{Precision of the economic comparison}
The target was a paired total width of $1/256$, fixed before seeing the new rows. Appendix~\ref{app:factorial81} reports every interval and its width components, including every marginal interval whether or not its width falls below that paired target. Variance reduction by conditional integration affects the variance component, not the bounded-support penalty or the acquisition contract. For example, in the fine calibrated basis comparison the conditional interval's width consists of approximately $0.001008605$ from the finite-support term, $0.000048152$ from endpoint sampling-variance allowances, and $2.69\times10^{-11}$ from the numerical enclosure. The support term is therefore the principal remaining contribution in that comparison.

The numerical component includes the implemented handling of continuous-bin width, acquisition ambiguity and arithmetic propagation. It is not represented as a pure floating-point error estimate. Likewise, a small numerical component does not justify dropping the finite-sample support term. These diagnostics explain the achieved precision without inferring significance from raw mean differences or treating the inherited R67 policy-cost rows as observations of these newly constructed actors.

\subsection{From a fitted readout to a qualified optimizer class}
For each new terminal-label qualification problem the action-switch catalogue covers every first-date observation cell and all thirty-two alternative robust actions. The 20,480 rational dual bounds are independently recomputed using the antiderivative and box-infimum calculations of Appendix~\ref{app:dual81}; no switch optimization is rerun to replace an unfavorable certificate.

\begin{table}[htbp]
\centering
\caption{First-date action stability of qualification sublevels}\label{tab:stability81}
{\small\setlength{\tabcolsep}{3.5pt}
\begin{tabular}{lrrrrc}
\toprule
Actor & Cells & $\epsilon=0$ & $1/65536$ & $1/4096$ & Uniform allowance \\
\midrule
CR8 & 64 & 28 & 28 & 28 & -- \\
CB8 & 64 & 64 & 64 & 55 & $1/32768$ \\
CR16 & 256 & 256 & 248 & 195 & $1/262144$ \\
CB16 & 256 & 256 & 243 & 200 & $1/8388608$ \\
\bottomrule\end{tabular}}
\par\smallskip{\footnotesize Counts use the strict inequality D greater than U plus epsilon. U is the exact objective of the attained feasible reference fit. The last column is a dyadic allowance strictly below the smallest positive switch surplus, when such a surplus is certified. A dash is unresolved, not evidence of instability. The independent verifier checks every one of the 20,480 action-switch duals.}
\end{table}

\begin{table}[htbp]
\centering
\caption{Cost comparisons uniform over exact qualification minimizers}\label{tab:optimizer-classes81}
{\small\setlength{\tabcolsep}{3.5pt}
\begin{tabular}{lc}
\toprule
Complete optimizer classes & Original-law cost difference \\
\midrule
CB16 minus CB8 & [-0.003286, -0.002241] \\
CB16 minus CR16 & [-0.002013, -0.000955] \\
CB8 minus CR16 & [0.000748, 0.001812] \\
\bottomrule\end{tabular}}
\par\smallskip{\footnotesize CB8, CR16 and CB16 have fully stable first-date covers at epsilon zero. Their terminal actors are fixed by the common terminal payoff. Consequently each stated exact-minimizer class gives the attained complete policy and inherits its interval without enlargement. This is a deterministic transfer of the same finite statistical family, not new sampling over optimizers.}
\end{table}


At zero extra objective allowance, the two fine-resolution constructions and coarse Bernstein have complete first-date action stability. Every exact minimizer of each of these fixed-label qualification problems therefore gives its attained complete policy, not merely a similar fitted value. The original-law intervals in Table~\ref{tab:optimizer-classes81} extend to these optimizer classes without additional probability expenditure. The same conclusion holds for the positive dyadic sublevels displayed in Table~\ref{tab:stability81}.

The standard $1/65536$ allowance leaves some fine-grid cells unresolved. Theorem~\ref{thm:stability81} then supplies an explicit enlargement based on their uniform initial mass. This conservative enlargement need not preserve a small cross-method sign. Coarse ReLU also retains unresolved cells even at zero allowance. None is relabeled as a discovered instability, and exact-minimizer invariance is not a claim that an arbitrary fitting routine reaches the requisite sublevel. The new statement addresses optimizer-independent policy stability under explicit attained-objective conditions; it does not estimate a hidden-layer success probability.

\subsection{Controller choice and complete costs}
Both cash and conditional marginal and pair intervals can be used together on their original simultaneous event. Exact shortest-path closure gives consistent cost potentials and the pure-choice regret bounds of the inherited decision theorem. The four-term interactions remain separately reported; sharpness of the pair-constraint calculation is not asserted after adding those extra restrictions. The complete nine-controller closure and attaining potentials are deposited in the revision audit.

For the fine calibrated controllers, the positive ReLU-minus-Bernstein interval determines a nonempty region of additional Bernstein charges for which its lower policy cost outweighs its extra charge. It also determines the opposite region in which ReLU is preferred after that charge. All such statements are conditional on an explicitly supplied conversion of computation, installation and deployment into normalized cost units. The exact boundary values may be tightened by the common pair-constraint closure; no new simulation is needed for that deterministic calculation.

The construction work is recorded rather than presumed equal. Legacy clocks include a label-loss ablation that the calibrated services do not execute. Their difference therefore does not identify an end-to-end speed gain from calibration. Exact action and expected-feature evaluation counts are reported alongside the clocks, and the original cost-inference, barrier and publication services retain separate measured costs. The present result supplies an identified policy-cost and conditional purchase comparison in the original economy. It does not extrapolate a two-date fitted-continuation calculation into a measured long-horizon speed result for the recursively queried vector tree.

\section{Economic stopping regions in the investment economy}\label{sec:study83}
The four calibrated constructions in Section~\ref{sec:study81} permit an exact application of the stopping theorem on $0\le\epsilon\le2^{-12}$. The full 20,480-switch catalogue leaves 976 actions for economic integration. All four declared dyadic levels are retained. Two independent arithmetic implementations reevaluate every paired box and reproduce the frozen records exactly. The tolerance calculation has 118 model strata; a separate enumeration checks all 12,732 pair-stratum tests. Appendix~\ref{app:frontier83} records this computation and the earlier interrupted execution separately.

At allowance $2^{-16}$, the fine-grid ReLU-minus-Bernstein class-cost band is $[0.000496,0.002503]$, with endpoints rounded outward. The corresponding binary-mass band is $[-0.421558,0.424526]$. Thus the same attained-policy evidence identifies a Bernstein saving throughout the approximately optimized classes once unresolved actions are priced individually. Every additional Bernstein deployment charge below $0.000496$ preserves the positive uniform net gain in normalized model units. This is not a calibrated installation fee.

\input{tables/frontier83}

The exact equal-tolerance boundary for that ordering is approximately $0.0000674917368$ and is excluded: a switch enters at equality. Its exact rational value, not the decimal approximation, governs a stopping rule. Cost equivalence within $1/256$ is a different question. Coarse Bernstein and fine ReLU satisfy this equivalence throughout $[0,2^{-12}]$, with final class-cost band $[-0.001622,0.003879]$. Both statements survive nonzero optimization error; neither is a universal representation ordering.

All four classes retain a positive cost saving against the zero controller throughout the integrated tolerance domain. Coarse ReLU nevertheless has inconclusive pairwise class comparisons already at zero excess. It and all opposite-ordering tests are not removed. The new frontier reuses the existing simultaneous cost event and supplies deterministic coverage for a continuum of verified stopping outputs. It creates neither new observations nor an estimate of fitting-success frequencies. Complete phase work, all discrete refinement bands, all eighteen frontier targets and both fixed-reference comparisons are retained in the appendix.

\section{Economic value of common-readout qualification}\label{sec:study84}
\input{tables/macros84}
The new calculation refines the economic accuracy of the original attained NBO controllers. It does not fit another controller or change a conditional cost observation. The four calibrated readouts, their frozen labels, terminal actors, robust 33-action menus, acquisition grids and original continuous laws are those of the preceding factorial investment study. The original actionwise integrations and simultaneous attained-policy intervals remain the paid inputs.

\subsection{A complete, fixed qualification design}
At acquisition sizes eight and sixteen, every alternative admitted by the reference-only barrier at excess objective at most $2^{-12}$ receives a full-menu qualification problem. This gives \FullMenuProblems\ problems: 832 for coarse ReLU, nine for coarse Bernstein, 78 for fine ReLU and 57 for fine Bernstein. Each event includes the other 32 feasible menu actions, with a common coefficient vector and the same objective.

The joint-event design permits at most 512 pairs per model. When at most 128 alternatives remain at the largest allowance, all remaining alternatives seed the pair catalogue. Otherwise the minimum-lower and maximum-upper conditional-cost alternatives in each cell seed it. Cross-cell pairs are ranked by the product of their maximum absolute cost endpoints, with lexicographic tie handling, and the first 512 are retained. The complete roster contains \JointProblems\ pairs; the small Bernstein model contributes all 36 of its possible pairs. This rule is deterministic postprocessing of an already covered cost event, not a newly randomized validation design. A small implementation pilot preceded the source lock and is disclosed separately.

All six declared allowances are evaluated: zero and $2^{-k}$ for $k=16,15,14,13,12$. The output preserves singleton bounds, pair bounds, positive box separations, retained alternatives, active conflicts and both economic duals. Every one of the \ExactInfeasible\ infeasible optimization statuses has an independent positive rational separation certificate. A numerical status alone is never used to remove an event.

\input{tables/compatibility84}

The full-menu test gives a strict contraction of the inherited interval in eight of the 24 model--allowance records. Shared-coefficient restrictions give a further strict contraction in fifteen records. Some smaller-tolerance records already have no alternative action or no useful selected conflict and remain unchanged. The table and appendix retain both outcomes. For coarse ReLU at the largest allowance, the inherited deviation interval \CoarseReLUBefore\ becomes \CoarseReLUAfter. This is a reduction in the uncertainty about a class of possible policies, not another measured gain of the attained reference policy.

\subsection{A new classwide refinement decision}
All six pairwise comparisons at the largest allowance appear in Table~\ref{tab:comparisons84}. They add the compatible deviation bounds to the same attained-cost intervals used in the preceding study. The result is simultaneous over the declared coefficient classes, rather than a comparison of two numerical certificate values.

\input{tables/comparisons84}

For coarse versus fine Bernstein, the previous optimizer-class interval at allowance $2^{-12}$ is \BernsteinResolutionBefore. It does not identify a signed ordering. The common-readout interval is \BernsteinResolutionAfter, whose lower endpoint is strictly positive. Thus the finer Bernstein controller has lower original-law cost for every pair of readouts in the two stated sublevels. Every additional normalized refinement charge below \BernsteinRefinementFee\ is supported by this classwide saving. This conclusion holds throughout smaller allowances by containment of the qualification classes, not by interpolation between the six printed tolerances.

The refinement decision is distinct from the matched basis comparison. At allowance $2^{-16}$ the fine ReLU-minus-Bernstein interval is \FineBasisAfter. It retains the identified Bernstein advantage under calibrated qualification while narrowing the preceding actionwise band. The earlier legacy qualification advantage of ReLU remains part of the factorial evidence. The paper therefore identifies the qualification rule and its verified optimization accuracy as economically relevant, rather than choosing a basis label and suppressing the opposite ranking.

The complete catalogue also retains coarse-ReLU comparisons that remain unresolved, and larger-allowance fine-basis comparisons whose intervals cross zero. An interval strictly inside the declared $1/256$ margin supports the corresponding equivalence statement; crossing zero alone does not. The new support calculation supplies no population fitting-success probability and does not convert a normalized cost unit into an estimated currency amount. The original priced-selection theorem can instead use the verified interval together with a separately supplied loss scale, deployment count and refinement charge.

\subsection{Complete incremental work and independent arithmetic}
The recorded scientific interpreter, including imports, reads, all qualification and assignment problems, durable outputs and shutdown, takes \CompleteScientificSeconds\ seconds on the stated ARM runner. Singleton, joint-event and assignment phases are recorded separately. These measurements price the additional qualification calculation. They are neither a cold-start cost for rebuilding the original controllers nor a matched ReLU-versus-Bernstein timing experiment. Original fitting and integration clocks remain unchanged and must be added when those inputs have not already been purchased.

A separate rational implementation reconstructs the feature antiderivatives, all winner comparisons, qualification supports, pair-event separations, conflict rosters and \IndependentInequalities\ assignment inequalities. It checks all 2,548 qualification certificates and all 36 policy-class comparisons without invoking the production optimizer. Original R77/R80 continuous-bin and actionwise kernel verification is separately replayed in the full publication build. The distinction between new exact certificate verification, inherited full-bin re-integration and independent cost sampling is explicit: this study adds no new fitting attempt, integration bin or independent cost row.

The economic improvement is consequently a more informative decision from the original NBO continuation and its previously qualified evidence. Full-menu winner regions remove alternatives that cannot be selected; common-coefficient restrictions remove jointly impossible selections; conditional integration prices what remains. Each step has its own proof, numerical certificate and work account. The original broader results and the difficult vector-horizon observations retain their separate assumptions and evidence.

\section{Prospective choice and economically interpretable comparisons}\label{sec:prospective86}
Exact ranking on a completed catalogue and preference before a new workload are different claims. We connect them by a population selection bound and a separately frozen evaluation, retaining the original investment economy and every comparator.

\begin{proposition}[Population regret for a fixed service menu]\label{prop:selection86}
Condition on construction of $m$ complete services, each with original-law cost in $[0,H]$. For $n$ independent new workloads, obtain valid pathwise cost enclosures $[l_{ji},u_{ji}]$. Write $\bar l_j,\bar u_j$ for endpoint means and $q_n=H\sqrt{\log(2m/\delta)/(2n)}$. Select the service minimizing $\lambda\bar u_j+p\cdot w_j+f_j$, with a fixed tie rule. Here $\lambda\ge0$, $p\ge0$, and the nonrandom resource and fee inputs are declared separately. With probability at least $1-\delta$, simultaneously for every such price-fee vector, its expected objective regret is at most
\[
 \lambda\{2q_n+\max_j(\bar u_j-\bar l_j)\}.
\]
A separate evaluation stream with simultaneous pairwise intervals $[l_{jk},u_{jk}]$ gives cost regret at most $\max\{0,\max_k u_{\widehat j,k}\}$ for any calibration-selected service $\widehat j$. This comparison is with the best expected-cost member of the declared menu, not the original Bellman optimum; the original-optimum account remains the separate policy sandwich.
\end{proposition}

\subsection{The frozen independent evaluation}
The prospective protocol preceding the retained evaluation fixes the nine complete two-date investment controllers: the zero controller and both ReLU and Bernstein readouts, under both legacy and calibrated qualification, at acquisition sizes eight and sixteen. Coefficients, complete action tables, constraints and original continuous innovation laws are unchanged. There are 8,192 calibration bins and 131,072 distinct evaluation bins, with forty bits per coordinate. Selection uses the smallest calibration conditioned upper-endpoint mean, with lexical ties, and is fixed before evaluation.

Both independently written whole-bin evaluators run on every policy and bin; rational midpoint fixtures check their hulls. The two estimators are cash cost and analytically conditioned cash cost. Each stream has 36 unordered pairs for each estimator. The empirical-Bernstein error account used in the original study assigns $4\exp(-11)$ to each interval, yielding $288\exp(-11)<0.005$ per stream and $576\exp(-11)<0.01$ jointly. No fitting-success event is inferred. The 139,264 paired workload rows count once, rather than once per policy or per comparison. Reconstructing the frozen streams during the current build adds no independent observations.

\begin{table}[htbp]
\centering
\caption{Independent evaluation: selected fine calibrated Bernstein minus comparator}\label{tab:selection86}
\begin{tabular}{lrr}\toprule
Comparator & Lower endpoint & Upper endpoint\\\midrule
Zero & $-0.177441$ & $-0.170788$\\
Legacy ReLU, size 8 & $-0.010137$ & $-0.005786$\\
Legacy Bernstein, size 8 & $-0.087292$ & $-0.081642$\\
Legacy ReLU, size 16 & $-0.009460$ & $-0.005046$\\
Legacy Bernstein, size 16 & $-0.086152$ & $-0.080530$\\
Calibrated ReLU, size 8 & $-0.009056$ & $-0.004812$\\
Calibrated Bernstein, size 8 & $-0.004823$ & $-0.000716$\\
Calibrated ReLU, size 16 & $-0.003559$ & $0.000573$\\\bottomrule
\end{tabular}
\par\smallskip
{\footnotesize Original-law conditioned cost. The table retains the prespecified R85 evaluation checkpoint; the current reproduction regenerates it from exact integer moments. Endpoints are rounded outwards to six decimal places. Selection uses 8,192 calibration rows, not the 131,072 evaluation rows. The simultaneous evaluation family contains 72 intervals across cash and conditioned estimators. The margin is $1/256$ in normalized cost units.}
\end{table}


Calibration selects the fine calibrated Bernstein controller. In the retained evaluation, its upper expected-cost regret is less than $0.000573$, below the prespecified $1/256$ noninferiority margin. Seven of its eight comparisons have negative upper endpoints. Its comparison with fine calibrated ReLU crosses zero and supports equivalence at that margin, not strict basis superiority. The resolution comparison within calibrated Bernstein is signed. The crossed basis-by-qualification effect is a separate factorial estimand, not a change in the underlying economy. Full endpoint moments, all 144 stream-estimator contrasts, sampling and enclosure components, and all comparisons are retained with the reproduction record.

The margin is a declared acceptable normalized loss. It is not an estimated monetary adoption fee. This evaluation does not replace the historical recursive-query workload, and does not assert that the nine acquired controllers are a prospectively validated neural query proposal on every new regime or horizon. The population construction in Section~\ref{sec:population86} supplies an explicit broader guarantee under its own joint-domain hypotheses; the two empirical objects are not conflated.

\subsection{Uniform price-fee sensitivity}\label{sec:prices86}
Let $\mathcal X$ be a finite set of feasible service assignments, including the original capacity constraints. Let $\mathcal D$ be a nonempty compact polytope containing their joint normalized cost vector. For an assignment $x$, let $w_x$ be its complete resource vector, $\lambda\ge0$ the value of normalized policy loss, $p\ge0$ the resource prices, and $f$ the assignment fee vector. All coordinates refer to explicit units. Put $h_{\mathcal D}(v)=\sup_{c\in\mathcal D}c\cdot v$.

\begin{proposition}[Uniform purchase region]\label{prop:prices86}
Assignment $x$ is optimal for every cost vector in $\mathcal D$ exactly on the region
\begin{equation}\label{eq:price-region86}
 \mathcal R_x=\bigcap_{y\in\mathcal X}
 \{(\lambda,p,f):\lambda h_{\mathcal D}(x-y)
              +p\cdot(w_x-w_y)+f\cdot(x-y)\le0\}.
\end{equation}
This region is polyhedral. Rational support-function upper certificates give an independently checkable sufficient region. Uniformity over a compact polytope of price-fee inputs can be checked at its vertices. The supporting equalities give the breakpoints at which a guaranteed choice can change. A point outside the sufficient region is not automatically evidence that the candidate is inferior.
\end{proposition}

For a single selected service and comparator $k$, a simultaneous cost upper bound $u_{\widehat j,k}$ gives the transparent sufficient condition
\[
 f_{\widehat j}-f_k+p\cdot(w_{\widehat j}-w_k)
 \le-\lambda u_{\widehat j,k}.
\]
The numerical appendix prints every selected-service breakpoint, including negative ones. In particular, the fine-versus-coarse calibrated Bernstein policy comparison supports a positive extra-fee allowance even at zero computation price, whereas a comparison whose upper cost endpoint is positive cannot justify a positive fee solely from policy quality. A resource-driven choice is distinguished from a policy-cost-driven one. Application-specific prices, energy costs and monetary scales require independent calibration; none is manufactured from the normalized tolerance or the favorable timing repetition.

\paragraph{Evidence identities.}
The inherited experiments have distinct observation units. The retained R67 description records 212 service returns and 48 neural fitting warnings; the later training, validation, reuse and vector catalogue records 188 complete processes. The attained-policy factorial study instead has nine controllers and its own independent path rows. The current complete-menu calculation has 2,548 qualification certificates and no new fits. Counts from different catalogues, including the 336-service description in the advisory report, are not added together or assigned to this evaluation without a source-matched manifest. The finer-bin experiment does not erase the older adverse timing and horizon results. Every inherited active section is retained, and the reply identifies the relevant experiment for each contested assertion.

\section{A frozen shared proposal on unseen workloads}\label{sec:sharedstudy89}
The joint-domain construction in Section~\ref{sec:population86} establishes a mathematical shared-operator guarantee. A trained proposal's incremental numerical value is a different question. We test that value in the original investment model without changing its continuous innovation law, feasible action cap or protected Bellman certificate. This experiment adds a prospective comparison to the inherited recursive-query and complete-policy studies; it does not substitute one of those objects for another.

\subsection{Construction, selection and evaluation}
The protocol is fixed before fitting. The training stream contains 128 states in each of dimensions two and eight, with two remaining dates. Protected adaptive search at tolerance $1/256$ supplies feasible action labels, normalized by the original action cap. Both approximation families receive the same fourteen features: eight zero-padded state coordinates, the overall, even-coordinate and odd-coordinate means, the mean square, the mean adjacent product, and the dimension divided by eight. The neural model is one width-32 ReLU network shared across both dimensions. The nonneural comparator is one complete quadratic on the same features. Each family has one attempted fit and fixed optimization settings, with no hyperparameter search. All coefficients, training rows and fitting diagnostics are retained before calibration.

Calibration has 64 new states in each training dimension. The five services are the adaptive baseline, bisection, a one-date analytic proposal, the shared ReLU proposal and the shared quadratic proposal. Each runs at tolerance $1/64$ using the unchanged protected oracle. Selection minimizes their total action-value queries, with the prespecified tie order. A terminal query is already included in that total and is not added a second time. Calibration selects the adaptive baseline: its 990 queries compare with 999 for bisection, 1008 for the analytic proposal, 996 for ReLU and 999 for the quadratic proposal. This choice is committed before the evaluation states are generated.

Evaluation contains 128 fresh states in each of dimensions two, four and eight. Dimension four is absent from both training and calibration. All services receive the identical states; three fresh-process repetitions preserve the complete timing distribution. The saved coefficients and evaluation inputs also define the cross-architecture replay. Repetitions and machine replay do not add independent states. All streams are reproducible seeded numerical workloads. A seed alone is not an assumption establishing population randomness; Proposition~\ref{prop:selection86} and the coverage theorem state the separate assumptions needed for probabilistic extrapolation.

\begin{table}[htbp]
\centering\caption{Shared proposals on the held-out original-law workload}\label{tab:work89}
{\small\setlength{\tabcolsep}{3pt}\begin{tabular}{lrrrrrr}
\hline
$d$ & Service & Nonterminal & Terminal & Model calls & Mean queries & Excess \\
\hline
2 & Adaptive & 325 & 650 & 0 & 7.61719 & 0 \\
2 & Bisection & 332 & 664 & 0 & 7.78125 & 21 \\
2 & Analytic & 333 & 666 & 0 & 7.80469 & 24 \\
2 & ReLU & 329 & 658 & 69 & 7.71094 & 12 \\
2 & Quadratic & 334 & 668 & 69 & 7.82812 & 27 \\
4 & Adaptive & 331 & 662 & 0 & 7.75781 & 0 \\
4 & Bisection & 333 & 666 & 0 & 7.80469 & 6 \\
4 & Analytic & 338 & 676 & 0 & 7.92188 & 21 \\
4 & ReLU & 339 & 678 & 75 & 7.94531 & 24 \\
4 & Quadratic & 334 & 668 & 75 & 7.82812 & 9 \\
8 & Adaptive & 323 & 646 & 0 & 7.57031 & 0 \\
8 & Bisection & 323 & 646 & 0 & 7.57031 & 0 \\
8 & Analytic & 325 & 650 & 0 & 7.61719 & 6 \\
8 & ReLU & 329 & 658 & 67 & 7.71094 & 18 \\
8 & Quadratic & 323 & 646 & 67 & 7.57031 & 0 \\
\hline
\end{tabular}}
\par\smallskip{\footnotesize Each row covers the same 128 states in its dimension. All services meet tolerance $1/64$. Mean queries counts nonterminal and terminal queries once. Excess is total queries above the best declared service in that held-out dimension, not population or policy-cost regret. Dimension four is absent from training and calibration. Three repetitions and the second machine do not multiply the sample size.}
\end{table}


\subsection{Results and their economic interpretation}
The two fitting attempts return finite coefficients without a recorded warning. All held-out services meet the common certificate tolerance. The adaptive service attains the smallest total query count in each evaluation dimension; at dimension eight it ties bisection and the quadratic proposal. Its empirical query regret against the declared menu is therefore zero in each dimension. This is not a zero expected-regret theorem and not a comparison with the original Bellman optimum.

The ReLU proposal uses 987, 1017 and 987 queries in dimensions two, four and eight, respectively, against 975, 993 and 969 for the adaptive service. The predictor is genuinely shared and is evaluated on unseen states and an unseen dimension, but it does not reduce verification work in this experiment. The protected first partition already resolves many workloads before any model call. In the remaining workloads, a model's proposed query can fail to improve the protected baseline's allocation. All such calls and refinements are charged. This result identifies the distinction between a valid common operator, a successfully returned fit and a cost-saving proposal; none implies the next.

Complete process times include interpreter startup, coefficient loading, computation, serialization and shutdown. The numerical appendix reports every repetition and peak resident memory on each recorded architecture, rather than only the fastest execution. Logical dense-model arithmetic and disjoint oracle queries remain separate accounts. They are not measured CPU instructions, energy expenditures or monetary costs. Since the calibration objective prices only action-value queries, its selected service need not minimize process time or a different priced resource vector. The price-region analysis in Section~\ref{sec:prices86} gives the appropriate conditional decision rule once those prices and complete resource accounts are specified.

\paragraph{Protection is not prediction success.}
The diagnostic catalogue separately records prediction-free returns, attempted screens, screen closures, protected refinements, invalid predictions, feasibility repairs and zero-progress probes. No failed fit or invalid prediction is omitted; the prescribed safe fallback remains a service outcome, not successful neural prediction. The two fits in this experiment do not erase the inherited nonconvergence warnings or estimate their population frequency. Original-law policy-cost conclusions continue to come from the separately frozen nine-controller evaluation in Section~\ref{sec:prospective86}, not from this query-proposal table.

\section{Deployment choice under qualified cost uncertainty}\label{sec:allocation79}
The attained continuation is an input to a complete economic decision, not merely to another approximation score. The same qualified costs can support a choice of acquisition precision across deployments. We retain the original NBO actors and draw a controller before observing the initial state or innovations. This implements a mixture of complete policies, not an average of their actions, coefficients or information sets. Its expected original-law cost is the corresponding mixture of the qualified costs. Minimax-regret allocation has an established decision-theoretic interpretation \citep{manski2004}; the result below concerns the finite cost set supplied by NBO qualification, conditional on its simultaneous coverage event.

Let $\mathcal C$ and $d(j,i)$ be the consistent difference-constraint set and distances of Theorem~\ref{thm:choice78}, stated and proved in the Numerical and Development Appendix. For a probability vector $p\in\Delta_q$ and per-deployment charges $w$, set $H_{ji}(w)=d(j,i)+w_i-w_j$, for $i,j=1,\ldots,q$. The charges include all costs that actually vary linearly with deployment shares; additional installation costs are treated separately.

\begin{theorem}[Sharp deployment allocation]\label{thm:allocation79}
The largest expected regret of allocating complete controllers with probabilities $p$ is
\begin{equation}\label{eq:allocation79}
 \mathcal R(p,w)=\max_{\theta\in\mathcal C}
 \left\{\sum_i p_i(\theta_i+w_i)-\min_j(\theta_j+w_j)\right\}
 =\max_j\sum_i p_iH_{ji}(w).
\end{equation}
Thus the optimal allocation and its value $V(w)$ are the solution of a finite linear program. For any rational $p,\lambda\in\Delta_q$, the computable numbers
\[
 L=\min_i\sum_j\lambda_jH_{ji}(w),\qquad
 U=\max_j\sum_i p_iH_{ji}(w)
\]
satisfy $L\le V(w)\le U$. There exist optimal primal and dual vectors with $L=U$. On the original coverage event, every implemented allocation has expected catalogue regret at most $U$, simultaneously for all $p,w$, including choices made after examining the qualified intervals. The worst-case bound in \eqref{eq:allocation79} is attained by one of the cost vectors already constructed in Theorem~\ref{thm:choice78}.
\end{theorem}

This is a decision calculation from the purchased observations, not a new experiment or an assertion that uncertainty reduction changes the true controller costs. Pure choices remain available as vertices of $\Delta_q$. Consequently $V(w)\le\min_i\Gamma_i(w)$, with strict improvement possible when no pure controller is uniformly best. The comparator is the best single qualified controller at the true cost vector. The original Bellman-optimum guarantees and the empirical comparisons between fitted representations remain separate.

\paragraph{An exact acquisition-fee region.}
Let $A$ and $B$ denote the fine and coarse ReLU controllers. The qualified gain is $J_B-J_A\in[l,u]$, where $0<l<u$. All other controllers are dominated by $B$ when their charges are equal. Suppose only $A$ has the incremental per-deployment charge $f$. For $l<f<u$, the fine-controller share and sharp expected regret are
\begin{equation}\label{eq:two-allocation79}
 p_A^*(f)=\frac{u-f}{u-l},\qquad
 V(f)=\frac{(f-l)(u-f)}{u-l}.
\end{equation}
For $f\le l$, use $A$ alone; for $f\ge u$, use $B$ alone. At the fee midpoint, allocation halves the best pure worst-case regret: $(u-l)/4$ rather than $(u-l)/2$. This is an exact decision bound, not a newly measured expected-cost reduction. In the retained experiment these bounds are approximately $0.0002405736$ and $0.0004811472$, respectively; Table~\ref{tab:allocation79} gives outward endpoints from the original rational intervals.

Maintaining both controllers need not be free. If a nontrivial allocation incurs an additional normalized cost $\kappa$, not incurred by either pure choice, its advantage survives precisely when
\begin{equation}\label{eq:setup-threshold79}
 \kappa<\frac{\min\{f-l,u-f\}^{2}}{u-l}.
\end{equation}
The full installed-subset account and a deterministic finite-cohort implementation are proved in Appendix~\ref{app:allocation79}. They charge every selected controller and do not turn shared experimental qualification into an unmeasured cold-start saving. For $m$ deployments and monetary loss scale $v>0$, the two-controller setup allowance is $mv$ times the right side. This gives an economically explicit frontier conditional on the actual fee, loss scale and deployment count, without inventing their calibration.

\section{Conclusion}\label{sec:conclusion65}
Neural Bellman Operators connect attained continuation construction to feasible decisions under the original law. Witnesses and residual qualification provide the original-optimum error account; affine calibration identifies a representation-dependent qualification charge. Actionwise cost transfer and exact tolerance frontiers then turn verified optimization accuracy into a sufficient economic stopping condition.

The investment comparison identifies gains from calibrated qualification in both bases and retains the reversal of their earlier ranking. Some cost orderings and equivalence conclusions hold throughout nonzero optimization sublevels without identical actions or new samples. Unresolved coarse classes remain explicit. These conclusions neither assume universal neural superiority nor establish unrestricted fitting reliability.

Conditional purchases and deployment allocations require separately supplied economic scales and charges. Original failed reuse, adverse query timings, vector descendant growth and broader applications retain their records and assumptions. The result is a complete constructive link from an attained approximation to a qualified controller and an economically interpretable accuracy requirement.

\paragraph{Shared coefficients and economic accuracy.}
A qualified output is one controller, not an arbitrary independent choice in every observation cell. Full-menu and shared-readout certificates respect this distinction and supply a constructive refinement of the original actionwise stopping regions. The assignment relaxation need not be an attainable optimizer policy, and every numerical upper bound is therefore identified as an enclosure rather than an observed worst-case loss. The verified common-coefficient restriction strengthens the economic use of the original NBO construction while preserving its original scope and all adverse evidence.

The complete-policy extension propagates changed actions through their own occupation flows. Joint-domain construction establishes uniform shared-operator guarantees under explicit primitive geometry, including endogenous contexts. Stored continuations and optimally allocated dated error budgets give a constructive horizon account distinct from the inherited descendant-tree experiment. Prospective selection and uniform price-fee regions then separate statistical policy quality from the cost of computation. Nonlinear preferences and strategic deviations retain their own valid comparison principles, rather than inheriting an unsupported interpretation of an expected-cost table.


\paragraph{The original law remains the benchmark.}
Pattern-consistent occupations and transport-qualified residuals provide complementary complete-policy guarantees. The former retains shared decisions across the contingent plan; the latter prices the action-value effect of numerical transition approximation, including singular laws. Both extend the original NBO construction and preserve its original-optimum interpretation. Their proofs identify every required regularity, qualification and arithmetic condition. The prospective evidence and coupled accuracy account connect those conditions to measured or explicitly priced resources without replacing unmeasured costs by assumed savings.

\bibliographystyle{ecta-fullname}
\bibliography{references,references78,references84,references90}
\end{document}
